\chapter{Literature Review}
\label{chapter:literature_review}

Generative Adversarial Networks (GANs) represent one of the most influential
families of deep generative models introduced in recent years for image generation tasks.
Originally proposed by Goodfellow et al.~\cite{background:gan},
GANs provide a framework for learning complex probability distributions
through an adversarial optimization process between two neural networks:
a generator and a discriminator.\\
Unlike likelihood-based generative approaches, where the objective is to
explicitly estimate the probability density function of the observed data,
GANs belong to the class of implicit generative models.
The generator learns a transformation from a simple latent distribution into
samples that approximate the unknown data distribution, without requiring an
explicit formulation of the underlying probability density.

\section{Standard GAN model}
\label{sec:gans}
The objective of the generator (G) is to produce synthetic observations that are
indistinguishable from real data, while the discriminator (D) attempts to identify
whether a given observation originates from the true data distribution or from
the generator.
GANs are often formalized as a minimax game, where the generator seeks to minimize the value function V while the discriminator
aims to maximize it:

\begin{equation}
    \label{eq:GAN minimax}
\min _G \max _D V(D, G)=\mathbb{E}_{\boldsymbol{x} \sim p_{\text {data }}(\boldsymbol{x})}[\log D(\boldsymbol{x})]+\mathbb{E}_{\boldsymbol{z} \sim p_{\boldsymbol{z}}(\boldsymbol{z})}[\log (1-D(G(\boldsymbol{z})))]
\end{equation}

Where $p_{\text {data }}$ is the unknown data distribution and $p_{\boldsymbol{z}}$ is the known latent distribution, usually a multivariate Gaussian or uniform distribution.
The discriminator outputs a value $D(x)$, 
which represents the probability $p$ that $x$ is a real sample.
Its goal is to correctly classify the real samples as real and generated samples as fake,
effectively maximizing the likelihood of the observed data.
To measure this, the cross-entropy loss is adopted, as it is common in deep learning,
expressed as $plog(1-p)$. 
For real samples, the true label $y$ equals 1, 
while for generated samples the label is 0.
With this setup, the discriminator's objective becomes:
\begin{equation}
    \label{eq:discriminator objective}
    \max _D V(D)=\mathbb{E}_{\boldsymbol{x} \sim p_{\text {data }}(\boldsymbol{x})}[\log D(\boldsymbol{x})]+\mathbb{E}_{\boldsymbol{z} \sim p_{\boldsymbol{z}}(\boldsymbol{z})}[\log (1-D(G(\boldsymbol{z})))]
\end{equation}

where $\mathbb{E}_{\boldsymbol{x} \sim p_{\text {data }}(\boldsymbol{x})}[\log D(\boldsymbol{x})]$ define how good the discriminator recognize real samples.
and $\mathbb{E}_{\boldsymbol{z} \sim p_{\boldsymbol{z}}(\boldsymbol{z})}[\log (1-D(G(\boldsymbol{z})))]$ define how good the discriminator recognizes generated samples.\\
The generator's objective is to produce distribution parameters that maximize $D(x)$, pushing the discriminator to classify them as real:
\begin{equation}
    \label{eq:generator objective}
\min _G V(G)=\mathbb{E}_{\boldsymbol{z} \sim p_{\boldsymbol{z}}(\boldsymbol{z})}[\log (1-D(G(\boldsymbol{z})))]
\end{equation}

Where $\mathbb{E}_{\boldsymbol{z} \sim p_{\boldsymbol{z}}(\boldsymbol{z})}[\log (1-D(G(\boldsymbol{z})))]$ define how good the generator is in fooling the discriminator.


\begin{figure}[H]
    \centering
    \includegraphics[width=1.0\textwidth]{intro_concepts/images/gan_architecture.png}
    \caption{Standard GAN architecture}
    \label{GAN_architecture}
\end{figure}

Moreover, as proved in \cite{background:gan}, for a fixed generator, the optimal discriminator can be derived analytically, obtaining the following expression:

\begin{equation}
D_G^*(x)
=
\frac{p_{data}(x)}
{p_{data}(x)+p_g(x)}.
\label{eq:optimal_discriminator}
\end{equation}

Where $p_g(x)$ is the distribution induced by the generator, equivalent to the distribution of the generated samples $G(z)$ when $z \sim p_z(z)$. 
Substituting this expression into the original minimax objective gives:

\begin{equation}
V(G,D_G^*)
=
- \log(4)
+
2 \cdot
JSD(p_{data}||p_g),
\label{eq:GAN_JS}
\end{equation}

where $JSD$ denotes the Jensen--Shannon divergence between the real and
generated distributions. Therefore, minimizing the GAN objective is equivalent to minimizing the
Jensen--Shannon divergence between the generated distribution and the true
data distribution, so when the generator is also optimal $p_g = p_{data}$ then the objective function reaches its minimum at $-\log(4)$.
Unfortunately, this remains a theoretical result 
since it is not possible, usually, to compute the optimal discriminator in practice,
and the training of GANs is often unstable and prone to issues such as mode collapse and vanishing gradients which will be discussed more in detail later on.\\


Despite the simplicity of the original formulation, GAN training suffers from
several fundamental difficulties.
The main issues include unstable optimization, vanishing gradients, mode
collapse, and sensitivity to hyperparameter choices.\\
These problems arise mainly because the discriminator and generator are
optimized simultaneously in a non-convex game.
Unlike classical optimization problems, where a unique minimum can often be
identified, GAN training seeks a Nash equilibrium between two competing
players.\\

Arjovsky and Bottou~\cite{background:gan_train_problems} demonstrated that the original GAN
objective can lead to severe optimization difficulties when the generated and
real distributions lie on low-dimensional manifolds.
In high-dimensional spaces, the supports of $p_{data}$ and $p_g$ may have
little or no overlap, causing the Jensen--Shannon divergence to become
constant.
When the discriminator becomes optimal too quickly, it can perfectly
distinguish real and generated samples. In this situation 
$D(G(z))\rightarrow 0$
and the gradient received by the generator becomes extremely small, providing little information for improving its distributional approximation capabilities.
This phenomenon is known as the \textit{vanishing gradient problem}.\\

Furthermore, GAN optimization may exhibit oscillatory behaviour rather than
convergence.
Mescheder et al.~\cite{background:gan_train_problems2} showed that simultaneous gradient-based
optimization does not necessarily converge to the desired Nash equilibrium,
even when the theoretical optimum exists.\\

These limitations become particularly relevant when GANs are applied to
financial time series. Unlike image datasets, financial observations are characterized by non-stationarity,
 regime changes, noisy stochastic dynamics and heavy tails distribution.
Therefore, instability in adversarial training may directly affect the quality
of generated financial scenarios and the reliability of downstream applications
such as risk estimation and derivative pricing.

Academic findings highlighted in \cite{background:gan_time_series} suggest that one of the principal failure modes of GANs 
arises when the generator converges to a parameter configuration that yields only a limited variety of outputs.
Existing research indicates that the occurrence of mode collapse is attributed to the establishment of a local 
equilibrium within the minimax optimization framework.
A further critical challenge involves the vanishing gradient problem,
which emerges when the discriminator becomes overly proficient,
causing the generator's gradients to approach zero and consequently impeding its learning process.\\


The original GAN framework learns the unconditional distribution
$p_{data}(x)$ and generates samples independently of any additional
information. However, many practical applications require the generation of
samples conditioned on external variables. In particular, in financial
forecasting the objective is not simply to generate realistic asset prices,
but rather to estimate the conditional distribution of future prices given
historical market information.\\
Conditional Generative Adversarial Networks (CGANs), introduced by Mirza and
Osindero~\cite{background:cgan}, extend the original GAN architecture by incorporating
additional information into both the generator and discriminator networks.
The conditioning variable can represent any relevant contextual information,
such as class labels, historical observations, market indicators, or previous
time-series values.\\
Let $c$ denote the conditioning variable. In the context of financial
forecasting, $c$ may represent the observed trajectory of an asset over a
historical window $c_t=(x_{t-k+1},x_{t-k+2},...,x_t)$ where $k$ denotes the length of the historical observation window.\\
The objective of a conditional generator is to learn the conditional
distribution $p(x|c)$ and generate samples that are consistent with the provided context $c$.\\

The generator receives both the latent noise vector $z$ and the conditioning
information $c$ while the discriminator evaluates whether a sample is consistent with the
conditioning information.
The conditional discriminator therefore does not only determine whether a
sample is realistic, but also whether it is compatible with the given
condition.
The CGAN objective function is given by:


\begin{equation}
    \label{eq:CGAN minimax}
    \min _G \max _D V(D, G)=\mathbb{E}_{\boldsymbol{x} \sim p_{\text {data }}(\boldsymbol{x})}[\log D(\boldsymbol{x} \mid \boldsymbol{c})]+\mathbb{E}_{\boldsymbol{z} \sim p_{\boldsymbol{z}}(\boldsymbol{z})}[\log (1-D(G(\boldsymbol{z} \mid \boldsymbol{c})))]
\end{equation}

A standard way for training GANs can be summarized by the following algorithm:

\begin{algorithm}[H]
\caption{Standard Conditional CGAN Training}
\label{alg: cgan}
\begin{algorithmic}[1]
\REQUIRE Initial discriminator parameters $\theta_D$, initial generator parameters $\theta_G$, condition matrix $C$, output data $X$.
\WHILE{$\theta_G$ has not converged}
    \STATE Sample real data $\{x^{(i)}\}_{i=1}^m \sim \mathbb{P}_{data}$ and context $\{c^{(i)}\}_{i=1}^m$.
    \STATE Sample latent noise $\{z^{(i)}\}_{i=1}^m \sim p_z$.
    \STATE $\tilde{x}^{(i)} \leftarrow G(z^{(i)}, c^{(i)})$ \COMMENT{Generate synthetic samples}
    
    \STATE \COMMENT{Update Discriminator: Maximize log(D(x)) + log(1-D(G(z)))}
    \STATE $L_{D} \leftarrow -\frac{1}{m} \sum_{i=1}^m \left[ \log D(x^{(i)}, c^{(i)}) + \log(1 - D(\tilde{x}^{(i)}, c  ^{(i)})) \right]$
    \STATE $\theta_D \leftarrow \theta_D - \alpha \nabla_{\theta_D} L_{D}$ \COMMENT{Gradient descent and backpropagation}
    
    \STATE Sample latent noise $\{z^{(i)}\}_{i=1}^m \sim p_z$.
    \STATE \COMMENT{Update Generator: Maximize log(D(G(z))) } 
    \STATE $L_{G} \leftarrow -\frac{1}{m} \sum_{i=1}^m \log D(G(z^{(i)}, c^{(i)}))$
    \STATE $\theta_G \leftarrow \theta_G - \alpha \nabla_{\theta_G} L_{G}$ \COMMENT{Gradient descent and backpropagation}
\ENDWHILE
\end{algorithmic}
\end{algorithm}


\section{Wasserstein GANs and other improvements}
The instability of the original GAN and (inherently the CGAN as well) formulation motivated the development of
alternative training objectives.
A major contribution was introduced by Arjovsky, Chintala and Bottou
through the Wasserstein GAN (WGAN) framework~\cite{background:wgan}.
The key idea behind WGAN is to replace the Jensen--Shannon divergence used in
the original GAN objective with the Wasserstein distance, also known as the
Earth Mover distance which measures the minimum cost required to transform one
probability distribution into another.

The transition from traditional divergence-based
Generative Adversarial Networks to those utilizing the Wasserstein metric 
represents a significant shift in the geometric treatment of probability 
measures within deep learning. 
The primal definition of the Wasserstein-1 distance 
for two probability distributions, $\mathbb{P}_r$ (real) and $\mathbb{P}_g$ (generated) is given 
by:
\begin{equation*}
    W(\mathbb{P}_r, \mathbb{P}_g) = \inf_{\pi \in \Pi(\mathbb{P}_r, \mathbb{P}_g)} \mathbb{E}_{(x, y) \sim \pi} [\|x - y\|]
\end{equation*}
where $\Pi(\mathbb{P}_r, \mathbb{P}_g)$ denotes the set of all joint distributions $\pi(x, y)$ with marginals $\mathbb{P}_r$ and $\mathbb{P}_g$.
While this formulation provides a physically intuitive measure of "work" required to transform one distribution into another, in fact this calculates the minimum "cost" of transporting mass to transform one distribution into the other.
It is computationally intractable in the context of neural networks because the infimum over the space of 
all possible transport plans $\pi$ cannot be efficiently optimized via backpropagation as suggested in \cite{background:wgan}.
Moreover, even if the two distributions are discretized through a binning strategy, an high computational cost is associated with the simulation of terminal values which define the empirical distribution and the following
discretization and evaluation of the Wasserstein distance.\\

To bridge this gap, CWGAN models rely on the \emph{Kantorovich-Rubinstein duality} \cite{intro:duality}, which transforms the constrained minimization over transport plans into a maximization over a set of 1-Lipschitz continuous functions.
The application of this duality is critical because it replaces the search for a transport plan with a search for a function $f$ (often called the Critic) that maximizes the difference between the expected values of the two distributions.\\


The main changes that characterize the CWGAN with respect to the CGAN are:
\begin{itemize}
    \item The loss function of a standard GAN is based on Jensen-Shannon divergence minimization while the CWGAN is based on the Wasserstein distance, which provides a more stable training process and mitigates issues such as mode collapse.
    \item The discriminator, which in Wasserstein GAN architecture is called critic, is designed to output real-valued scores instead of probabilities, allowing it to better capture the differences between real and generated data distributions.
    For this reason, the discriminator network does not use a sigmoid activation function in the output layer.
    \item The use of gradient penalty term in the loss function has the objective to enforce the Lipschitz constraint required for the Wasserstein distance computation.    
\end{itemize}

\begin{theorem}[Kantorovich duality]
\label{eq:kantarovich}
Let $(\mathcal{X}, \mu)$ and $(\mathcal{Y}, \nu)$ be two Polish probability spaces and let $c : \mathcal{X} \times \mathcal{Y} \to \mathbb{R} \cup \{+\infty\}$ be a lower semicontinuous cost function, such that
\begin{equation*}
    \forall(x,y) \in \mathcal{X} \times \mathcal{Y}, \quad c(x,y) \geq a(x) + b(y)
\end{equation*}
for some real-valued upper semicontinuous functions $a \in L^1(\mu)$ and $b \in L^1(\nu)$. Then
\begin{itemize}
    \item[(i)] There is duality:
\end{itemize}
\begin{align*}
    \min_{\pi \in \Pi(\mu, \nu)} & \int_{\mathcal{X} \times \mathcal{Y}} c(x, y) \, d\pi(x, y) \\
    &= \sup_{(\psi, \phi) \in C_b(\mathcal{X}) \times C_b(\mathcal{Y}); \,\, \phi-\psi \leq c} \left( \int_{\mathcal{Y}} \phi(y) \, d\nu(y) - \int_{\mathcal{X}} \psi(x) \, d\mu(x) \right) \\
    &= \sup_{(\psi, \phi) \in L^1(\mu) \times L^1(\nu); \,\, \phi-\psi \leq c} \left( \int_{\mathcal{Y}} \phi(y) \, d\nu(y) - \int_{\mathcal{X}} \psi(x) \, d\mu(x) \right) \\
    &= \sup_{\psi \in L^1(\mu)} \left( \int_{\mathcal{Y}} \psi^c(y) \, d\nu(y) - \int_{\mathcal{X}} \psi(x) \, d\mu(x) \right) \\
    &= \sup_{\phi \in L^1(\nu)} \left( \int_{\mathcal{Y}} \phi(y) \, d\nu(y) - \int_{\mathcal{X}} \phi^c(x) \, d\mu(x) \right),
\end{align*}
and in the above suprema one might as well impose that $\psi$ be $c$-convex and $\phi$ $c$-concave.
\end{theorem}

To formalize the components of the Kantorovich duality
within the context of Wasserstein GANs (WGANs), 
we define the underlying mathematical structures and their functional roles.
In this framework, the pairs $(\mathcal{X}, \mu)$ and $(\mathcal{Y}, \nu)$ 
represent Polish probability spaces, where in the specific case of generative modeling,
$\mu$ corresponds to the real data distribution $P_r$ and $\nu$ represents the model's 
generated distribution $P_g$. The cost function $c(x, y) \colon \mathcal{X} \times \mathcal{Y} \to \mathbb{R} \cup \{+\infty\}$ serves as the metric 
for the "work" required to transport a unit of mass from a point $x$ in the real domain to a point $y$ 
in the generated domain. The central object of the primal problem is the transport plan $\pi \in \Pi(\mu, \nu)$, 
which is a joint distribution whose marginals are constrained to $\mu$ and $\nu$. 
This plan describes the optimal allocation of mass across the product space.
Transitioning to the dual formulation introduces the dual potentials $\phi$ and $\psi$.
In the architecture of a WGAN, the potential $\phi$ (often denoted as $f$) is parameterized as the Critic network.
These functions assign scalar values to points in the manifold, effectively measuring the "potential energy" difference used to calculate the Wasserstein distance. Mathematically,
these potentials reside in $L^1(\mu)$, the space of functions that are integrable with respect 
to the measure $\mu$. A critical component for achieving the final loss function is the $c$-transform (or $c$-conjugate) of $\phi$, defined as:

\begin{equation}
    \phi^c(x) = \inf_{y \in \mathcal{Y}} [c(x, y) - \phi(y)]
\end{equation}

This transform identifies the optimal "mate"
function that satisfies the feasibility constraint
$\phi(y) - \psi(x) \leq c(x, y)$. When the cost function $c(x, y)$
is chosen as the $L_1$ distance $\|x - y\|$, the $c$-transform forces 
$\phi$ to satisfy a Lipschitz continuity constraint. 
This transformation is what allows the WGAN to replace the intractable search over all transport plans $\pi$ with a simpler maximization over $1$-Lipschitz functions, 
providing the stable gradients necessary for training deep generative models.

So finally, the Kantorovich-Rubinstein duality can be expressed in the specific case of the Wasserstein-1 distance, given
$P_r$ and $P_g$ two probability distributions defined on a compact metric space $\mathcal{X}$ and where
the supremum is taken over all $1$-Lipschitz continuous functions $f: \mathcal{X} \to \mathbb{R}$, such that for all $x, y \in \mathcal{X}:   |f(x) - f(y)| \leq \|x - y\|$, then:
\begin{equation}
    \label{eq: emd_dual}
    W(P_r, P_g) = \sup_{\|f\|_L \leq 1} \left( \mathbb{E}_{x \sim P_r}[f(x)] - \mathbb{E}_{x \sim P_g}[f(x)] \right)
\end{equation}


In this duality, the "discriminator" is redefined as a "critic" ($D$), whose task is not to classify
inputs as real or fake, but to act as the optimal function $f$ that maximizes the distance between the 
expected scores of the real and synthetic distributions.
In the following thesis, $f$ is parametrized by using a neural network $f_w$ with weights $w$ restricted to a compact space $\mathcal{W}$ to enforce the Lipschitz constraint.
Substituting the generator output $G(z)$ for samples from $P_g$, the optimization objective function for the Wasserstein GAN is obtained.
The minimax objective function is the following:
\begin{equation}
    \label{eq:wgan_objective_function}
    L(G, f_w) = \min_G \max_{w \in \mathcal{W}} \mathbb{E}_{x \sim P_r}[f_w(x)] - \mathbb{E}_{z \sim p(z)}[f_w(G(z))]
\end{equation}
By enforcing $\|f_w\|_L \leq K$ (where $K$ is a constant, typically through weight clipping or gradient penalty), the Critic provides a linear, non-saturating gradient to the Generator, ensuring stable training even when the Generator is far from the target distribution.

In a Conditional GAN (CGAN), this distance is 
extended to account for auxiliary information $c$ (such as past history, drift, or volatility in 
time-series). So in the conditional setting, 
mathematically, the EMD is calculated over the pairs $(x, c)$, where $c$ is the history and $x$ is the generated output.
The real pair $(x_{i, real}, c_i)$ and fake pair $(G(z_i, c_i), c_i)$ are produced. The "Earth mass" is not being moved on a 1D line but instead it is being moved in a higher-dimensional space where one axis is the context and the other is the generated output. 
The distance is only minimized when the generator produces the correct $x$ for that specific $c$.

The objective becomes the minimization of the Wasserstein distance between the 
conditional distributions $\mathbb{P}_r(x|c)$ and $\mathbb{P}_g(x|c)$. 
By conditioning both the generator $G(z, c)$ and the critic $D(x, c)$ on the same context $c$,
the critic learns a "score surface" that is specific to the underlying stochastic process defined by $c$.

The use of the critic's output to compute the loss is mathematically 
justified because, under the 1-Lipschitz constraint, 
the difference between the average scores of the real batch and 
the fake batch directly approximates the Wasserstein distance as demonstrated in \cite{background:wgan}.
This is a superior training signal compared to Jensen-Shannon divergence because
$W(\mathbb{P}_r, \mathbb{P}_g)$ is
continuous and differentiable even when the supports of the two
distributions are disjoint, a common occurrence in the early stages of training.
Consequently, the generator receives stable, non-vanishing gradients that indicate 
not just if the generated sample is wrong,
but how far it must move in the feature space to align with the real
data manifold. 
To enforce the Lipschitz constraint required by the duality, modern implementations utilize a Gradient Penalty (GP),
which regularizes the norm of the critic's gradient with respect to its input toward unity, ensuring the function remains 
within the space of 1-Lipschitz functions $\mathcal{D}$.

\begin{algorithm}[H]
\caption{Conditional Wasserstein GAN with Gradient Penalty (CWGAN-GP)}
\label{alg:cwgan_gp}
\begin{algorithmic}[1]
\REQUIRE The gradient penalty coefficient $\lambda$, the number of critic iterations per generator iteration $n_{critic}$, the batch size $m$, Adam hyperparameters $\alpha, \beta_1, \beta_2$.
\REQUIRE Initial critic parameters $\theta_D$, initial generator parameters $\theta_G$.
\WHILE{$\theta_G$ has not converged}
    \FOR{$t = 1$ \TO $n_{critic}$}
        \STATE Sample real data $\{x^{(i)}\}_{i=1}^m \sim \mathbb{P}_r$ and context $\{c^{(i)}\}_{i=1}^m$.
        \STATE Sample latent noise $\{z^{(i)}\}_{i=1}^m \sim \mathbb{P}_z$.
        \STATE $\tilde{x}^{(i)} \leftarrow G(z^{(i)}, c ^{(i)})$ \COMMENT{Generate synthetic samples}
        \STATE $\epsilon \sim U[0, 1]$ \COMMENT{Random number for interpolation}
        \STATE $\hat{x}^{(i)} \leftarrow \epsilon x^{(i)} + (1 - \epsilon) \tilde{x}^{(i)}$ \COMMENT{Interpolated samples} 
        \STATE $L_{D} \leftarrow \frac{1}{m} \sum_{i=1}^m D(\tilde{x}^{(i)}, c^{(i)}) - \frac{1}{m} \sum_{i=1}^m D(x^{(i)}, c^{(i)}) + \lambda \mathbb{E}_{\hat{x}, c} [(\|\nabla_{\hat{x}} D(\hat{x}, c)\|_2 - 1)^2]$  
        \STATE $\theta_D \leftarrow \text{Adam}(\nabla_{\theta_D} L_{D}, \theta_D, \alpha, \beta_1, \beta_2)$
    \ENDFOR
    
    \STATE Sample latent noise $\{z^{(i)}\}_{i=1}^m \sim \mathbb{P}_z$.
    \STATE $L_{G} \leftarrow -\frac{1}{m} \sum_{i=1}^m D(G(z^{(i)}, c^{(i)}))$
    \STATE $\theta_G \leftarrow \text{Adam}(\nabla_{\theta_G} L_{G}, \theta_G, \alpha, \beta_1, \beta_2)$
\ENDWHILE
\end{algorithmic}
\end{algorithm}

An import part of the training procedure is the gradient penalty calculation, which is described in Algorithm \ref{alg:gradient_penalty}.
By sampling $\epsilon$ and creating $\hat{x}$, the Lipschitz constraint is enforced specifically on the straight lines between the real and fake data manifolds.

\begin{algorithm}[H]
\caption{Gradient Penalty Calculation for CWGAN-GP}
\label{alg:gradient_penalty}
\begin{algorithmic}[1]
\REQUIRE Real samples $x$, generated samples $\tilde{x}$, conditioning context $c$, critic network $D$.
\STATE $m \leftarrow \text{batch size}$
\STATE Sample $\epsilon \sim U[0, 1]^m$ \COMMENT{Uniform random interpolation weights}
\STATE $\hat{x} \leftarrow \epsilon x + (1 - \epsilon) \tilde{x}$ \COMMENT{Compute interpolated samples}
\STATE Enable gradient tracking for $\hat{x}$
\STATE $\hat{S} \leftarrow D(\hat{x}, c)$ \COMMENT{Evaluate critic on interpolates}
\STATE $\nabla_{\hat{x}} D \leftarrow \nabla_{\hat{x}} \hat{S}$ \COMMENT{Compute gradients of scores w.r.t. interpolates}
\STATE $L_2 \text{-norm} \leftarrow \|\nabla_{\hat{x}} D(\hat{x}, c)\|_2$ \COMMENT{Calculate norm across feature dimensions}
\STATE $GP \leftarrow \frac{1}{m} \sum (L_2 \text{-norm} - 1)^2$ \COMMENT{Compute mean squared error from unity}
\RETURN $GP$
\end{algorithmic}
\end{algorithm}

So the final objective function for the Conditional Wasserstein GAN with gradient penalty (CWGAN-GP) becomes:

\begin{equation}
    \begin{aligned}
        \min _G \max _D V(D, G) = & \mathbb{E}_{\boldsymbol{x} \sim p_{\text {d}}}[D(\boldsymbol{x} \mid \boldsymbol{c})] - \mathbb{E}_{\boldsymbol{z} \sim p_{\boldsymbol{z}}}[D(G(\boldsymbol{z} \mid \boldsymbol{c}))] + \lambda GP
    \end{aligned}
\end{equation}
where $\lambda$ is a hyperparameter that controls the strength of the penalty and it is usually set to $10$ as suggested in \cite{background2:gp_wgan}.\\
An additional advantage of using the Wasserstein distance is that it provides a meaningful measure of the distance between the real and generated data distributions, which allows for the implementation of an early stopping technique which
improve the training efficiency.\\

When the generated distribution and real distribution are far apart, the
Jensen--Shannon divergence often becomes constant.
Consequently, the discriminator saturates and the generator receives
vanishing gradients. Conversely, the Wasserstein distance changes continuously as the generated
distribution moves toward the target distribution.
Therefore, even when the two distributions have limited overlap, the critic
provides meaningful gradient information.\\

Other tentatives to improve the stability of GAN training. One of the most successful techniques are the spectral normalization and feature matching.
Miyato et al.~\cite{background:spectral_norm} proposed Spectral Normalization as a method to
control the Lipschitz constant of neural networks. The weight matrix of each discriminator layer is normalized according to its
spectral norm: $\bar{W}=\frac{W}{\sigma(W)}$
where $\sigma(W)$ represents the largest singular value of $W$.
Spectral normalization constrains discriminator updates and improves
stability without requiring additional gradient penalties.
It has become one of the most effective stabilization techniques for modern
GAN architectures.\\

Salimans et al.~\cite{background:feature_matching} introduced feature matching as a method to
reduce mode collapse. Instead of directly maximizing discriminator classification error, the
generator minimizes the difference between statistics of intermediate
discriminator features:

\begin{equation*}
L_G
=
\left|
\mathbb{E}_{x\sim p_{data}}D(x)
-
\mathbb{E}_{z\sim p_z}D(G(z))
\right|^2 .
\end{equation*}

This encourages the generator to reproduce multiple modes of the target
distribution rather than exploiting weaknesses of the discriminator.


\section{GANs for Time Series}
The architectures reviewed so far in this chapter, 
the original GAN, its conditional extension, and the variance-reduction/stabilization techniques, were developed 
primarily with static, i.i.d. data, most prominently images, in mind. 
Extending the adversarial framework to \emph{sequential} data introduces a distinct
set of architectural requirements: 
the generator and discriminator must be able to represent
temporal dependence, like autocorrelation, rather than merely a single, 
unconditional multivariate distribution, and, for variable-length or streaming applications, 
must do so with a fixed-size internal state rather than a fixed-size input/output tensor. 
This section reviews the principal families of sequential GAN architectures, 
culminating in the financial-time-series-specific 
proposals that directly motivate the modelling choices made later in this thesis.\\

The most direct way to extend a GAN to sequential data is to replace the feed-forward generator and discriminator networks of the original architecture with recurrent networks (LSTMs or GRUs) that process a sequence one step at a time,
propagating a hidden state $\mathbf{h}_\tau$ forward across the sequence index $\tau$. 
A generic recurrent generator produces a sequence $\hat{\mathbf{x}}_{1:H} = (\hat x_1, \dots, \hat x_H)$ autoregressively,
\begin{equation*}
\hat x_\tau = g_\theta\left(\mathbf{h}_{\tau-1}, z_\tau\right), \qquad \mathbf{h}_\tau = f_\theta\left(\mathbf{h}_{\tau-1}, \hat x_\tau, z_\tau\right), \qquad \tau = 1,\dots,H,
\end{equation*}

\cite{background:rcgan} propose the Recurrent Conditional GAN (RCGAN), developed originally for the generation of realistic multivariate medical time series. RCGAN uses an RNN (LSTM) for both generator and discriminator, with the generator receiving \emph{independent} noise draws at every time step rather than a single noise vector at initialization,
\begin{equation*}
\hat x_\tau = G_\theta\left(\mathbf{h}_{\tau-1}, z_\tau\right), \qquad z_\tau \stackrel{\text{iid}}{\sim} \mathcal{N}(0, I),
\end{equation*}
a design choice the authors show empirically increases sample diversity relative to single-noise-vector variants, directly counteracting the temporal mode-collapse tendency whereby a generator conditioned on a single static noise draw tends to produce sequences that are smooth, low-entropy interpolations rather than genuinely stochastic paths. The discriminator likewise processes the full sequence recurrently and is trained with the standard (or, in later variants, Wasserstein) adversarial loss applied either to the final hidden state or averaged across per-step outputs. RCGAN additionally supports conditioning on auxiliary static or time-varying labels.\\

Also \cite{background:timegan} propose a novel architecture for time series generation that addresses a specific shortcoming of purely adversarial recurrent GANs: 
the adversarial loss alone provides no explicit incentive for the 
generator to preserve the \emph{stepwise conditional} dynamics 
of the training data (i.e. the transition structure $p(x_\tau \mid x_{1:\tau-1})$), since the discriminator, 
in its simplest form, only ever evaluates entire generated sequences 
(or fixed-length windows) against real ones in aggregate, without an explicit term that penalizes mismatches in 
the one-step-ahead transition law. TimeGAN addresses this by combining the adversarial objective with a 
\emph{supervised} autoregressive loss, and by decoupling the 
space in which the generator operates from the original data space via a 
learned embedding.\\

The proposed architectures for time series proposed so far are designed for the generation of realistic sequences, but do not directly address the problem of \emph{forecasting}, 
where the goal is to estimate the conditional distribution of future values given a history. One of the most influential architectures in this space is ForGAN \cite{background:forgan}, 
which combines the conditional adversarial framework with a recurrent conditioning network.

\section{Probabilistic Forecasting using GANs}
The overwhelming majority of forecasting models developed in the quantitative finance and machine learning literature are trained to minimize a pointwise loss function, typically the mean squared error (MSE) or the mean absolute error (MAE), between a predicted scalar or vector and an observed realization. Formally, given a conditioning history $\mathbf{x}_{1:t} = (x_1, \dots, x_t)$ of a stochastic process, a point forecaster learns a deterministic mapping
\begin{equation*}
\hat{x}_{t+h} = f_\theta(\mathbf{x}_{1:t}), \qquad \theta^\star = \arg\min_\theta \; \mathbb{E}\left[ \ell\left(x_{t+h}, f_\theta(\mathbf{x}_{1:t})\right) \right],
\end{equation*}
where $\ell(\cdot,\cdot)$ is typically the squared error $\ell(x,\hat x) = (x-\hat x)^2$.
When the underlying data-generating process is stochastic, as is the case for essentially every financial time series, since asset prices, volatilities, and interest rates evolve under genuine aleatoric uncertainty, a point forecast conveys only the conditional mean (or, depending on the loss, the conditional median) of the future state, and discards all information about the shape, skewness, and tail behaviour of the conditional distribution.
This is a critical shortcoming for financial applications, where risk management, derivatives pricing, and portfolio construction depend fundamentally on the full conditional law of future returns rather than on a single expected value.\\

The correct object of interest is therefore the conditional predictive distribution
\begin{equation}
p\left(x_{t+h} \mid \mathbf{x}_{1:t}\right),
\end{equation}
or, in the multi-step-ahead case relevant to path-dependent risk metrics, the joint conditional distribution
\begin{equation}
p\left(x_{t+h}, x_{t+2}, \dots, x_{t+H} \mid \mathbf{x}_{1:t}\right)
\end{equation}
over a forecast horizon $H$. 
Estimating this object is the task of \emph{probabilistic forecasting} \cite{background:probabilistic_forecasting}. 
Classical approaches to probabilistic forecasting in econometrics rely on explicit parametric assumptions, for example, GARCH-family models coupled with a Gaussian or Student-$t$ innovation distribution,
or the Heston stochastic volatility model itself, 
which specifies a closed-form (up to Fourier inversion) 
characteristic function for the conditional log-price distribution. 
These approaches are interpretable and computationally tractable, 
but their accuracy is contingent on the correctness of the parametric assumption; 
misspecification of the innovation distribution or of the volatility dynamics propagates directly into biased risk estimates.\\

Generative Adversarial Networks 
offer an attractive alternative precisely
because they are implicit generative models
where the generator $G(\mathbf{c}, \mathbf{z})$, conditioned on a history $\mathbf{c} = \mathbf{x}_{1:t}$ 
and driven by an exogenous noise vector $\mathbf{z} \sim p_z$, is trained so that its induced output distribution approximates $p(x_{t+h} \mid \mathbf{x}_{1:t})$
without any parametric assumption on the shape of that distribution. 
Repeated sampling of $\mathbf{z}$ for a fixed conditioning history $\mathbf{c}$ produces an empirical ensemble of trajectories, from which arbitrary functionals of the predictive distribution, quantiles, moments, tail probabilities and so on,
can be estimated by Monte Carlo. This is the conceptual bridge between adversarial generative modelling and probabilistic time series forecasting,
and it is the foundation of the ForGAN architecture discussed below.\\

ForGAN, introduced by \cite{background:forgan},
is a conditional GAN architecture designed
explicitly for one-step-ahead probabilistic
forecasting of univariate time series. 
It combines the conditional adversarial training framework
with a recurrent conditioning network, so that the generator 
can encode an arbitrarily long history into a fixed-dimensional context vector 
before producing a sample of the next-step distribution.\\

The ForGAN's architecture is based on a generator and a discriminator, both utilizing Recurrent Neural Network (RNN) layers, specifically LSTM or GRU variants, in order to process sequential data.\\
The generator's objective is to forecast a future next value $x_{t+h}$ by transforming the noise vector $z$, 
conditioned on the historical window $c$.\\
The probability distribution of $x_{t+h}$ is modeled as $P(x_{t+h}|c)$ and obtained by running the generator multiple times with different noise inputs and same condition window $c$.\\
Structurally, the generator processes the condition window $c$ through an RNN layer to construct a "condition representation" whose size is given by hidden state size.
This representation is subsequently concatenated with the noise vector $z$.
The combined feature vector is then passed through two dense layers, culminating in a single scalar output representing the predicted $x_{t+h}$.\\
On the other hand, the discriminator has the role to inspect the validity of a given value $x_{t+h}$ in the context of the history $c$.
Unlike standard discriminators that might process inputs separately, the ForGAN discriminator takes the candidate $x_{t+h}$, which is sourced either from the ground truth dataset or the generator, and concatenates it to the end of the condition window $c$, forming an extended sequence $\{x_{0},...,x_{t+h}\}$.
This extended sequence is fed into an RNN layer, followed by a dense layer that outputs a validity score, determining whether the sequence is real or generated.
So the innovation is given by the presence of a conditioner which is implemented as a recurrent neural network that processes the conditioning window $\mathbf{c} = (x_{t-\ell+1}, \dots, x_t)$ of length $\ell$ and returns the terminal hidden state as a fixed-length embedding,
\begin{equation}
C_\phi(\mathbf{c}) = \mathrm{RNN}_\phi(x_{t-\ell+1}, \dots, x_t) \in \mathbb{R}^{d_c}.
\end{equation}\\

\begin{figure}[H]
    \centering
    \includegraphics[width=1.0\textwidth]{Introduction/images/forgan_rnn.png}
    \caption{ForGAN building blocks}
    \label{img:forgan_rnn}
\end{figure}

\begin{figure}[H]
    \centering
    \includegraphics[width=1.0\textwidth]{Introduction/images/forgan.png}
    \caption{ForGAN architecture}
    \label{img:forgan}
\end{figure}

The model is trained with the standard conditional minimax objective adapted to the time series conditioning structure. In fact, as proposed in \cite{background:forgan}, the objective is to minimize the Jensen-Shannon divergence between the real and generated data as seen before \ref{eq:CGAN minimax}.
Because the original ForGAN objective inherits the training pathologies
of the vanilla GAN, like vanishing gradients when the discriminator saturates,
and susceptibility to mode collapse (see Section~\ref{sec:limitations-gans}), 
a natural extension replaces the Jensen--Shannon-based minimax loss with the Wasserstein critic loss \ref{eq:wgan_objective_function}.
In this last case the model can be referenced as ForWGAN which represent a novel 
approach to increase the performance of the standard ForGAN.\\

While ForGAN 
was shown to perform strongly on the sensory
and synthetic benchmarks
for which it was originally developed, \cite{forgan:fin_gan} directly investigate 
its behaviour on financial return data, 
and show that a ForGAN trained with its original, 
generic binary-cross-entropy (BCE) adversarial loss is \emph{systematically outperformed},
on economically meaningful metrics, by both classical econometric baselines (ARIMA)
and by a variant of the same ForGAN architecture trained with a purpose-built, 
economics-driven loss function, the proposed Fin-GAN model.
This result is important for the present thesis because it demonstrates 
that the limitations catalogued in Section~\ref{sec:limitations-gans} 
are not merely a theoretical curiosity: 
they have a directly measurable, negative impact on downstream financial performance,
specifically in the ForGAN architecture that the SRForGAN model developed in this thesis builds upon.\\
The core diagnosis in \cite{forgan:fin_gan} 
is a form of \emph{objective misalignment} between what the ForGAN 
adversarial loss actually optimizes and what a 
financial forecasting application actually requires. 
In the Fin-GAN paper, the authors solve the inherent misalignment of the BCE-trained ForGAN by replacing the generic adversarial loss with a composite loss function that explicitly encodes 
the financial decision-making objectives of interest, 
including mean-squared error, sign accuracy, profit-and-loss, and risk-adjusted return. 
So, they worked on the loss function to make it more aligned with the financial forecasting task, 
rather than on the architecture itself. It is indeed, the same approach that is pursued in this thesis.\\


The Fin-GAN critique of BCE-trained ForGAN is,
in an important sense, structurally analogous to the motivation 
for the scoring-rule-based training paradigm underlying the SRForGAN architecture developed in this thesis.
Both approaches share the diagnosis that the original ForGAN training objective,
precisely because it is a generic adversarial divergence with no reference 
to the specific quantity the downstream application 
cares about, is an unreliable foundation for financial forecasting,
and both approaches respond by replacing that generic objective with one explicitly aligned 
to a well-defined target: 
in Fin-GAN's case, an economics-driven, 
decision-rule-specific composite loss 
(MSE, sign accuracy, PnL, and risk-adjusted-return terms); 
in the case of this thesis, a strictly proper scoring rule whose global minimizer is, 
by construction, the true conditional \emph{distribution} itself, 
rather than a decision-rule-specific functional of it.

\section{Scoring Rules and ForGAN advances}



%-------------------------------------------------------------
%-------------------------------------------------------------




\label{sec:energy_score_theory}
A scoring rule is any function $S:\mathcal P\times\Omega\to\mathbb R$ 
that assigns a numerical reward $S(P^\phi,y)$ to a 
probabilistic forecast $P\in\mathcal P$ once the 
outcome $y$ materialises \cite{background:scoring_rules}. 
Scoring rules originated as tools for \emph{forecast evaluation} and 
\emph{elicitation}. A forecaster who 
believes the true law is $P^\star$ is rewarded, 
in expectation, according to $S(P^\phi,P^\star)$ for
 whichever forecast $P^\phi$ is quoted, and 
 propriety is precisely the condition 
 guaranteeing that quoting $P^\phi =P^\star$ is the 
 forecaster's optimal, honest strategy.  
 Beyond evaluation, any strictly proper $S$ 
 can be repurposed as a \emph{training objective}: 
 minimising $-S(P^\phi, P^\star)$ over a parametric or 
 generative family $\{P^\phi\}$ pushes $P^\phi$ 
 toward the data-generating law, which is the 
 perspective adopted throughout this thesis.\\

This section formalises why minimising the empirical 
Energy Score of Eq.~\eqref{eq:energy_score_estimator} 
drives the generator$G_\phi$ toward the 
true conditional law $P^\star(\cdot|c)$, 
and states the regularity conditions under 
which this guarantee holds for the i.i.d.\ 
trajectory design of the subsequent experiments discussed in the following chapter.\\

Propriety is the most important charactersistic 
that a scoring rule can have and it is extensively leveraged during this work.
The authors of \cite{background:scoring_rules} show that propriety 
 can be fully characterized:
 \cite{background:scoring_rules} 
 show that every proper
 scoring rule derives
from a convex function $F$ on 
$\mathcal P$, where $F(P):=S(P,P)$ 
is the expected score under the forecast $P$ itself. 
Propriety is equivalent to $F$ being convex and 
to $S(P,\cdot)$ acting as a subtangent of 
$F$ at $P$; strict propriety corresponds 
to $F$ being strictly convex. $F$ 
thus plays the role of a generalized 
entropy, and the divergence $D(P^\star\|P^\phi)$ 
becomes, on suitable spaces, the 
Bregman divergence induced by $F$. 
Any admissible convex $F$ therefore 
yields a propriety-preserving training loss that can be optimized by gradient descent,
which is precisely what allows scoring rules 
to be tailored to a given problem.

\paragraph{Strict propriety.} 
For a scoring rule $S(P^\phi,y)$, 
define the expected score 
$S(P^\phi,P^\star):=\mathbb{E}_{Y\sim P^\star}S(P^\phi,Y)$. 
Following \cite{background:scoring_rules},
 $S$ is \emph{proper} relative to a class $\mathcal{P}$ 
 if $S(P^\star,P^\star)\le S(P^\phi,P^\star)\ \forall P^\phi\in\mathcal{P}$, 
 and \emph{strictly proper} if 
 equality holds only at $P^\phi=P^\star$. 
 Propriety is what makes $S$ a sound training 
 objective: it guarantees
$$D(P^\star\|P^\phi):=S(P^\phi,P^\star)-S(P^\star,P^\star)\ \ge\ 0,$$
is a valid statistical divergence, vanishing
iff $P^\phi=P^\star$.\\

The authors of \cite{background:scoring_rules} 
uses this property to demonstrate that given a fitting problem
of a parametric model $P^\phi$ where $Y_1, Y_2, \ldots, Y_n$ are the observed i.i.d. 
samples, and $\phi_0$ the true parameter value, then by 
maximizing the goodness-of-fit given by the expected score:

\begin{equation*}
    \mathcal{S}_n(\phi)=\frac{1}{n} \sum_{i=1}^n S\left(P_\phi, Y_i\right),
\end{equation*} 

then: $\arg \max _\phi \mathcal{S}_n(\phi) \rightarrow \phi_0 \text { as } n \rightarrow \infty$.\\


%Any $\phi$ minimising $\mathbb{E}_{c\sim\Pi}\,S(P^\phi(\cdot|c),P^\star(\cdot|c))$ therefore recovers the \emph{full}  conditional law for every $c\in\mathrm{supp}(\Pi)$,  not merely a point estimate or a moment of it.
 
One of the most widely used strictly proper scoring rules for multivariate continuous distributions is the Energy Score, which generalises the Continuous Ranked Probability Score (CRPS), which is the main scoring rule used in this work, to higher dimensions.\\
$S_E^{(\beta)}$ is defined as:
\begin{equation}
\mathrm{S}_E^{(\beta)}(P^{\phi}, \mathbf{y}) = \mathbb{E}\left\| \mathbf{X} - \mathbf{y} \right\|_d^\beta - \tfrac{1}{2}\, \mathbb{E}\left\| \mathbf{X} - \mathbf{X}' \right\|_d^\beta, \qquad \mathbf{X}, \mathbf{X}' \stackrel{\text{iid}}{\sim} P^{\phi},
\end{equation}

which is basically the Kernel Score $S_k(P^\phi,y)=\mathbb{E}[k(X,X')]-2\mathbb{E}[k(X,y)]$ 
with the negative-power kernel $k(x,y)=-\|x-y\|_2^\beta$. 
In \cite[Thm.~4]{background:scoring_rules}, authors show that
$S_k$ is strictly proper on the class of distributions with 
finite $\mathbb{E}\|X\|^\beta$ whenever $k$ is (conditionally) 
strictly positive-definite. For $\beta\in(0,2)$ the negative-power kernel satisfies this,
so $S_E^{(\beta)}$ is strictly proper on that class and by consequence also the univariate CRPS which is obtained by setting $\beta=1$.\\ 


%----------------------------------------------
Before using $S_E^{(\beta)}$ as a training loss, \cite{srforgan:srforgan} and \cite{background:scoring_rules} require:
 
\begin{itemize}
\label{scoring_rule_assumptions}
\item $\beta\in(0,2)$: the negative-power kernel $k(x,y)=-\|x-y\|^\beta$ is (conditionally) strictly positive-definite only in this range, which is what makes $S_E^{(\beta)}$ strictly proper rather than merely proper.
\item Finite $\beta$-th moments: $\mathbb{E}_{Y\sim P^\star}\|Y\|^\beta<\infty$ and $\mathbb{E}_{X\sim P^\phi}\|X\|^\beta<\infty$, without this the score, and hence the objective, may not even be finite.
\item A \emph{well-specified} model class: some $\phi^\star\in\Phi$ with $P^{\phi^\star}(\cdot|c)=P^\star(\cdot|c)$ for all $c$, needed for the population minimiser to be unique (\cite{srforgan:srforgan}, Lemma~12). This condition is satisfied by the theorem of the universal approximation of neural networks, 
which guarantees that a sufficiently large generator network can approximate any continuous function to arbitrary precision.
\item i.i.d.\ training pairs $(c_j,y_j)\sim\Pi\times P^\star(\cdot|c)$: if pairs are instead dependent (e.g.\ sliding windows of one sequence), consistency additionally requires stationarity and mixing conditions (\cite{srforgan:srforgan}, Thm.~3), which our design avoids by construction.
\item A generator admitting the reparameterisation trick: $X=G_\phi(z,c)$, $z\sim\mathcal N(0,I)$, differentiable in $\phi$, so $\nabla_\phi S_E$ can be estimated unbiasedly (App.~C, \cite{srforgan:srforgan}).
\end{itemize}



For our i.i.d.\ simulation 
design, this is not a modelling risk but an analytically verifiable property: 
with $\beta=1$, the log-price increment over $[X_{Ndt},X_{Ndt+\tau}]$ is Gaussian
 with variance $\sigma_j^2\Delta t$ and $\sigma_j\in(0,1.2]$ bounded above by 
 construction (Algorithm~\ref{alg:y condition}), so $\mathbb{E}|Y|<\infty$ trivially, 
 and $\mathbb{E}\|X\|<\infty$ follows for any generator with bounded output layer or 
 empirically bounded activations. This is a strictly easier position than the 
 prequential time-series case of \cite{srforgan:srforgan}, where the same 
 moment condition must hold \emph{uniformly over $t$} under stationarity and 
 mixing assumptions (their Theorem 3), an assumption we do not need, since our 
 $J$ trajectories are simulated i.i.d.\ rather than drawn as overlapping windows 
 of one sequence.\\

The proper scoring rules most relevant to this literature and in particular to this work is the Continuous Ranked Probability Score (CRPS) which is derived from the Energy Score by using a beta parameter equal to 1.
For a univariate predictive CDF $\hat F$ and a scalar outcome $y$:
\begin{equation}
    \label{eq:CRPS}
\mathrm{CRPS}(\hat F, y) = \int_{-\infty}^{\infty} \left( \hat F(x) - \mathbf{1}_{x \ge y} \right)^2 \, dx,
\end{equation}
which admits the equivalent, sample-based (kernel) representation in terms of two independent draws $X, X' \sim \hat F$,
\begin{equation}
\mathrm{CRPS}(\hat F, y) = \mathbb{E}\left| X - y \right| - \tfrac{1}{2}\, \mathbb{E}\left| X - X' \right|.
\end{equation}
This representation is directly usable as a training loss for a GAN generator: given $N$ samples $\{\hat x^{(i)}\}_{i=1}^N$ drawn from $G([\mathbf{c},\mathbf{z}^{(i)}])$, an unbiased estimator of the CRPS is
\begin{equation}
    \label{eq:CRPS_estimator}
\widehat{\mathrm{CRPS}} = \frac{1}{N} \sum_{i=1}^{N} \left| \hat x^{(i)} - x_{t+h} \right| - \frac{1}{2 N (N-1)} \sum_{i=1}^{N} \sum_{j \neq i} \left| \hat x^{(i)} - \hat x^{(j)} \right|,
\end{equation}


The Energy Score, and by consequence the CRPS, therefore plays a role directly analogous to the minimax value function $V(D,G)$ of the original GAN, but with the crucial advantage that it is a proper scoring rule with a known, well-behaved minimizer, removing the discriminator-induced gradient pathologies discussed in Section~\ref{sec:limitations-gans} at the cost of a bias--variance trade-off governed by the ensemble size $N$ (the estimator above is unbiased for the population Energy Score, but its variance decreases only at rate $O(1/N)$, so small ensembles yield noisy gradients.


\section{Applications in Quantitative Finance}
 
The appeal of adversarially or scoring-rule
trained generative models in quantitative 
finance stems from a structural feature 
shared by nearly every downstream task in the field: 
risk management, derivative pricing, and portfolio 
optimization all reduce, at some level, 
to computing an expectation (or a quantile, or a full distribution) 
of a functional of future asset paths. 
Classical Monte Carlo methods evaluate 
such expectations by simulating paths 
from an explicit, calibrated stochastic model 
(e.g. Heston, SABR, or a local-volatility model); 
GAN-based methods replace 
this explicit simulator with a 
learned, data-driven generator, 
potentially side-stepping model 
misspecification at the cost of 
requiring large amounts of 
representative historical data and 
forfeiting the interpretability of a closed-form model.\\

Distribution-valued
forecasting has a growing dedicated literature in finance.
\cite{background:quantile_regression} use a deep quantile-regression
network to forecast Value-at-Risk directly, 
reporting significant accuracy gains over linear quantile regression. 
A recent study \cite{background:applications} 
trains CNN and LSTM networks to output the parameters 
of skewed-$t$ predictive distributions, 
evaluating them with CRPS and PIT diagnostics 
and showing they are competitive with GARCH for 
VaR estimation. 
\cite{background:applications2} combine deep quantile regression with GAN-based scenario generation specifically to estimate VaR and Expected Shortfall jointly, explicitly motivating GAN-generated tail scenarios as a remedy for the difficulty classical models have capturing stressed regimes.\\

On the pricing side, generative and simulation-based deep models have been applied to option-market 
simulation and hedging, 
notably \cite{background:applications3}'s Deep Hedging framework and 
 Quant GAN for scenario generation. 
Also the work of \cite{background:applications4}, Tail-GAN, modify the adversarial objective itself to 
target VaR/ES elicitability directly, 
reflecting the same insight, 
that generic distributional losses under-serve the tail, 
that motivates using a proper scoring rule rather than a raw adversarial loss in SRForGAN. 
These works establish both the demand for, and the 
empirical viability of, GAN and neural-network-based 
probabilistic forecasters as inputs to standard risk and 
pricing pipelines, and SRForGAN can be positioned within 
this literature representing a suitable GAN-based architecture
capable of capturing the full conditional distribution of financial processes,
representing a valid alternative to parametric models and other generative architectures.


\section{Limitations of standard forecast techniques}
\label{sec:literature_gap}

Classical forecasting casts $x_{t+h}$ prediction as a regression problem,
with the model trained to minimize a pointwise loss such as RMSE. 
Under squared-error loss, the Bayes-optimal predictor is the conditional expectation,
\begin{equation}
\label{eq:mean_regression_target}
\hat{x}_{t+h} = \arg\min_{f} \; \mathbb{E}_{x_{t+h}}\big[(x_{t+h} - f(c))^2\big] = \mathbb{E}[x_{t+h} \mid c] = \mu(\rho(x_{t+h}\mid c)),
\end{equation}
i.e. the model converges
to the mean of the conditional distribution 
regardless of its shape. 
This is only informative when $\rho(x_{t+h}\mid c)$ 
is unimodal and concentrated, approximating a Dirac delta. 
When the true conditional distribution is multimodal, Eq.~\ref{eq:mean_regression_target} 
produces a value that need not lie near \emph{any} 
realizable outcome. 
\cite{background:forgan} illustrate this with a synthetic bimodal target where $x_{10}$ 
realizes one value in $80\%$ of cases and a second, 
disjoint value in the remaining $20\%$: 
the trained regressor converges to the 
weighted average $0.8$, 
a point with zero support 
under the true data-generating process. 
This failure mode is structural rather 
than a symptom of under-fitting or 
insufficient capacity. 
No amount of additional training data or 
model complexity corrects it, since the objective
in Eq.~\ref{eq:mean_regression_target} is itself misspecified 
relative to the estimand of interest 
whenever $\rho(x_{t+h}\mid c)$ deviates 
from near-degeneracy. 
The same argument applies to the Heston-driven volatility clusters targeted in this thesis, 
where the conditional density of log-price increments is heteroskedastic 
and can exhibit heavy tails and skewness 
inconsistent with a single-mode Gaussian-like summary.

Quantile regression \cite{background:quantile_regression} partially 
addresses this limitation by estimating 
conditional quantile functions $Q_{\alpha}(x_{t+h}\mid c)$ directly, 
via the asymmetric piecewise-linear (pinball) loss
\begin{equation}
\label{eq:pinball_loss}
\mathcal{L}_{\alpha}(x_{t+h}, \hat{x}_{t+h}) = \big(x_{t+h} - \hat{x}_{t+h}\big)\big(\alpha - \mathbf{1}[x_{t+h} < \hat{x}_{t+h}]\big), \qquad \alpha \in (0,1),
\end{equation}

consistent for the $\alpha$-quantile. 
Fitting a family $\{\hat{Q}_{\alpha_1}, \dots, \hat{Q}_{\alpha_K}\}$ over a grid of 
quantile levels in principle recovers an approximation 
of the full conditional distribution. 
In practice, however, each $\hat{Q}_{\alpha_k}$ 
is typically estimated by an independent regression, 
with no constraint enforcing monotonicity in $\alpha$. 
This gives rise to the \emph{quantile-crossing} problem: 
for $\alpha_i < \alpha_j$ the fitted quantiles 
can satisfy $\hat{Q}_{\alpha_i}(c) > \hat{Q}_{\alpha_j}(c)$, 
violating the definition of a valid distribution function. 
The pathology is most pronounced with dense quantile grids or 
small sample sizes \cite{background:mean_regression}, 
precisely the regime relevant to derivative-pricing 
and tail-risk applications where fine resolution in 
the extreme quantiles (VaR, ES) is required. Post-hoc 
rearrangement or monotonicity-constrained estimators 
\cite{background:quantile_regression2,background:quantile_regression3} 
mitigate but do not eliminate the problem, and add 
estimation overhead that scales with the number of 
quantile levels considered.

Both failure modes share a common root: 
they optimize a functional of $\rho(x_{t+h}\mid c)$ 
(its mean, or a finite set of quantiles) 
rather than the distribution itself, 
and the functional's fidelity to the 
underlying density degrades precisely in 
the multimodal, heteroskedastic regimes characteristic 
of stochastic volatility models. 
This motivates the shift toward implicit, 
sample-based conditional density estimation 
pursued by ForGAN and, subsequently, 
by the SRForGAN architecture proposed in this thesis.




\section{Limitations of GANs}
\label{sec:limitations-gans}
Despite their theoretical appeal as flexible, likelihood-free generative models, GANs are notoriously difficult to train in practice, and the literature documenting and diagnosing their failure modes is nearly as extensive as the literature proposing new architectures. This section reviews the principal sources of instability, with particular attention to the manifestations that are most consequential for financial time-series applications.\\

The original GAN formulation establishes tha, in the nonparametric limit, that is when both $G$ and $D$ have unlimited capacity and are optimized exactly, alternately, to their respective optima at each step,
the minimax game converges to a unique Nash equilibrium at which $p_g = p_{\mathrm{data}}$. This guarantee, however, does not extend to the practically relevant setting in which $G$ and $D$ 
are finite-capacity neural networks trained by simultaneous or alternating stochastic gradient descent. 
In this setting, the minimax objective \ref{eq:GAN minimax}
defines a non-convex, two-player game,
for which simultaneous gradient descent is not guaranteed to converge even locally,
let alone to the desired Nash equilibrium. 
The resulting optimization landscape is known to exhibit oscillatory, non-convergent trajectories even on simple, low-dimensional toy problems,
where the generator and discriminator can enter a persistent limit cycle rather than converging to the equilibrium.\\

One of the most widely documented failure modes is \emph{mode collapse}, 
refers to the failure mode in which the generator $G$ learns to produce outputs concentrated on a small subset of the modes
(or, in the extreme case, a single point or a very low-dimensional manifold) of the true data distribution
$p_{\mathrm{data}}$,
rather than covering its full support.
The phenomenon arises, as discussed in \cite{background:mode_collapse}, because the minimax objective, when optimized via alternating gradient steps
rather than the idealized best-response dynamic assumed in the original convergence proof,
provides no direct incentive for the generator to maintain output diversity:
if a subset of the noise space $\mathcal{Z}$ maps to outputs that reliably fool the current discriminator,
gradient descent on parameters of $G$ has no mechanism to \emph{penalize} 
the generator for ignoring the remainder of $\mathcal{Z}$, since the objective is evaluated only in expectation over $\mathbf{z}$ and contains no explicit diversity or entropy term.\\
Actually, gradient descent effectively optimizes the cases in which the generator 
need only to find, for each fixed discriminator, \emph{some} output that fools it, rather than a full distribution that matches $p_{\mathrm{data}}$
under every possible discriminator. 
\cite{background:wgan} further connect mode collapse to a distributional-support argument:
when $p_{\mathrm{data}}$ and $p_g$ have disjoint or low-dimensional-overlapping supports
(a generic situation when both are supported on low-dimensional manifolds embedded in a high-dimensional ambient space, as is typically the case for both images and, arguably, financial processes),
the Jensen--Shannon divergence between them is \emph{constant} (equal to $\log 2$)
and its gradient with respect to the generator's parameters is zero almost everywhere, 
providing no useful training signal and creating an optimization landscape in which the generator can become ``stuck'' 
having covered only a subset of the target modes, 
with no local gradient information indicating the direction of the missing modes.\\
In the context of probabilistic financial forecasting, 
mode collapse manifests specifically as \emph{underdispersion}: 
a generator conditioned on a given history $\mathbf{x}_{1:t}$ 
collapses across noise draws $\mathbf{z}$ 
to a narrow cluster of outputs, so that the induced 
predictive distribution $\hat p_G(x_{t+h} \mid \mathbf{x}_{1:t})$ 
has substantially smaller variance than the true conditional distribution.\\

A second, related instability is the \emph{gradient vanishing} that arises from gradient saturation at the discriminator.
Recall the original generator loss $\mathbb{E}_{\mathbf z}[\log(1 - D(G(\mathbf z)))]$:
when the discriminator is properly trained, 
it rapidly learns to assign $D(x) \approx 1$
for real samples and $D(G(\mathbf z)) \approx 0$ 
for generated samples that are not yet realistic. 
In this regime,
\begin{equation}
\frac{\partial}{\partial \theta} \, \mathbb{E}_{\mathbf z}\left[ \log\left(1 - D_\psi(G_\theta(\mathbf z))\right) \right] \; \longrightarrow \; 0,
\end{equation}
with $\theta$ the discriminator parameters and $\phi$ the generator parameters. This happens because
the function $\log(1-y)$ has near-zero slope as $y \to 0$, 
so that precisely when the generator is performing poorly and most needs a strong corrective gradient,
the gradient it receives vanishes. 
\cite{background:gan} proposed the standard workaround, training the generator to instead maximize $\mathbb{E}_{\mathbf z}[\log D_\psi(G_\theta(\mathbf z))]$, which restores a non-vanishing gradient in the early-training, 
poorly-performing regime, but at the cost of abandoning the interpretation of the generator's objective 
as the minimization of a well-defined statistical divergence between $p_g$ and $p_{\mathrm{data}}$,
and reintroducing exactly the mode-seeking, 
diversity-agnostic behaviour discussed above. 
\cite{background:wgan} show more generally that, 
whenever the discriminator is allowed to approach optimality 
(which is precisely what is needed for it to provide an accurate estimate of any $f$-divergence, including the Jensen--Shannon divergence underlying the original GAN), 
the generator's gradient with respect to the Jensen--Shannon-based objective vanishes in the limit, 
establishing vanishing gradients not as an incidental implementation 
artifact but as a structural property of divergence-minimizing GAN objectives whenever $p_g$ and $p_{\mathrm{data}}$ 
have (near-)disjoint support, precisely the motivation for the Wasserstein reformulation discussed, since the Earth-Mover distance, unlike the Jensen--Shannon divergence, 
remains a smooth, non-vanishing function of the generator's parameters even under disjoint supports.\\

All these pathologies are the exact reason why the introduction
of scoring-rule-based GANs is so important for financial forecasting applications and 
motivate the development of the SRForGAN architecture in this thesis.
































\chapter{Data, Models and Experimental Setup}
Experiments are conducted to evaluate the performance of the proposed SRForGAN model in comparison to the original ForGAN architecture in order to assess their ability to produce coherent conditional distributions given well known financial processes, 
in particular, the experiments adopt trajectories obtained from Geometric Brownian Motion and Heston stochastic volatility model. 
This chapter outlines how training data are computed, model configurations, training procedures, and evaluation metrics used in the experiments.

\section{Training Data}
In order to effectively train the Conditional GAN models, two kinds of data are required for the training phase:
\begin{itemize}
    \item Condition input data $C$: it is a $J \times N$ matrix, where $J$ is the number of trajectories and $N$ is the number of time steps of each trajectory that goes from time 0 up to $T$.
     This data represents the past information available to the model to infer a coherent output of the asset price.
     \item Target output data $X$: Since in this thesis we focus on a single horizon output, this data is a vector of size $J$, where $J$ represents the number of trajectories.
    So, each entry of the vector contains the future output at time $T + \tau$ associated with the respective past trajectory. 
    \end{itemize}

However, in the test phase, a third kind of data is required, that is the ground truth data distribution $D_\tau(c)$ which is the 
distribution of the future output $\tau$ steps ahead given the past trajectory $c$. In this experimental setting, this distribution can be 
computed analytically in the case of the Black-Scholes model and 
numerically in the case of the Heston model (or alternatively using the inverse fast Fourier transform to obtain the probability density function as shown in \ref{alg:heston_pdf}).\\

\section{The ForGAN Architecture}
\label{sec:forgan}

ForGAN \cite{background:forgan} reformulates the $h$-step-ahead time series forecasting as a conditional generative modeling problem,
replacing the point estimate $\hat{x}_{t+h}$ of classical regression with an implicit sampler for the full conditional density $\rho(x_{t+h} \mid c)$, where $c = \{x_0, \dots, x_t\}$ denotes the historical condition window. 
This design directly targets the two structural weaknesses of 
mean-regression forecasters discussed in previous chapter \ref{chapter:literature_review}: convergence to the conditional mean under multimodal targets, and the quantile-crossing pathology of explicit quantile regressors.

ForGAN builds on the conditional GAN (CGAN) framework of \cite{background:cgan}, 
in fact, the mimised loss function is the same of a standard conditional GAN model:
\begin{equation*}
\min_G \max_D V(D,G) = \mathbb{E}_{x_{t+h} \sim \rho_{\text{data}}(x_{t+h})}[\log D(x_{t+h} \mid c)] + \mathbb{E}_{z \sim \rho_z(z)}[\log(1 - D(G(z \mid c)))].
\end{equation*}
Where the target data to generate is no more a matrix representing an image but the next value of a sequential process.
At the Nash equilibrium of this game, the optimal generator $G^*$ implicitly represents the true conditional distribution 
$\rho(x_{t+h}\mid c)$, so that repeated sampling $z \sim \mathcal{N}(0, I)$ yields draws from the forecast distribution rather than a single summary statistic.

Architecturally, both $G$ and $D$ are built around a single recurrent layer (LSTM or GRU). 
The generator first encodes the condition window through the RNN layer to 
obtain a fixed-size condition representation, 
concatenates it with a noise vector 
$z \sim \mathcal{N}(0, I)$ of size $N$, 
and passes the result through two dense 
layers to produce the scalar forecast $\hat{x}_{t+h}$ as represented in schemas \ref{img:forgan} and \ref{img:forgan_rnn}. 
The discriminator appends the candidate value (either $\hat{x}_{t+h}$ from $G$ or the ground-truth $x_{t+h}$) 
to the end of the condition window, 
forming the extended sequence $\{x_0,\dots,x_{t+h}\}$, 
and feeds it through its own RNN layer followed by a dense 
layer that outputs a single real/fake validity score. 
Training proceeds asymmetrically: at each epoch the discriminator 
is updated for $D_{\text{iter}}$ steps while the generator
 is updated once, a ratio tuned alongside the RNN cell type, 
 hidden sizes $R_G, R_D$, noise dimension $N$, and condition length $C$ via hyperparameter optimization techniques like genetic algorithms. 
 Algorithm~\ref{alg:forgan_training} summarizes one training epoch.

\begin{algorithm}[H]
\caption{ForGAN Training (single epoch)}
\label{alg:forgan_training}
\begin{algorithmic}[1]
\REQUIRE Generator $G_\phi$, discriminator $D_\theta$, condition window $c$, target $x_{t+h}$, discriminator iterations $D_{\text{iter}}$
\FOR{$i = 1$ to $D_{\text{iter}}$}
\STATE Sample noise $z \sim \mathcal{N}(0, I)$
\STATE $\hat{x}_{t+h} \leftarrow G_\phi(z, c)$ \COMMENT{generated forecast}
\STATE $\ell_D \leftarrow -\big[\log D_\theta(x_{t+h} \mid c) + \log(1 - D_\theta(\hat{x}_{t+h} \mid c))\big]$
\STATE Update $\theta$ by descending $\nabla_\theta \ell_D$
\ENDFOR
\STATE Sample noise $z \sim \mathcal{N}(0, I)$
\STATE $\hat{x}_{t+h} \leftarrow G_\phi(z, c)$
\STATE $\ell_G \leftarrow -\log D_\theta(\hat{x}_{t+h} \mid c)$
\STATE Update $\phi$ by descending $\nabla_\phi \ell_G$
\end{algorithmic}
\end{algorithm}
As can be denoted from the following algorithm is that 
it reproduce the same logic of the standard conditional GAN training process like in
\ref{alg: cgan}. In fact, an easy extension of the ForGAN model using the Wasserstein Loss
can be easily deduced by adopting the algorithm of the Wasserstein GAN with Gradient Penalty \ref{alg:cwgan_gp}, which is denominated ForWGAN during this thesis.\\

To isolate the contribution of adversarial training from that of the architecture itself, 
the original authors also define a \emph{G-regression} baseline: a network 
structurally identical to $G$, 
but trained by minimizing RMSE directly instead of playing the minimax game. 
Across the Lorenz, Mackey-Glass, and Internet-traffic (A5M) benchmarks, 
ForGAN is shown to match or beat G-regression and prior state-of-the-art 
on point-wise metrics (RMSE, MAE, MAPE) 
while dramatically outperforming both on the 
Kullback-Leibler divergence between the sampled 
forecast distribution and the empirical target distribution.\\
Despite this success, 
vanilla adversarial training inherits well-known 
instabilities: vanishing gradients when $D$ saturates, 
and mode collapse, 
whereby $G$ learns to reproduce only a subset of the 
modes of $\rho(x_{t+h}\mid c)$ while still fooling $D$ as described in section \ref{sec:limitations-gans}. 
These pathologies motivate the design choices underlying SRForGAN, 
introduced in the following section, 
which departs from the adversarial signal of Eq.~\ref{eq:CGAN minimax} 
in favor of a scoring-rule-based objective, 
retaining ForGAN's conditional generative structure 
while removing its reliance on a saturating discriminator game.


\section{The SRForGAN Architecture}
Mode collapse, also known as the "Helvetica scenario," 
occurs when the generator learns to produce only a 
limited subset of the data distribution as already explained in section \ref{sec:limitations-gans}. 
It can be categorized into two forms, especially in the context of 
conditional time-series GANs like ForGAN:
\begin{itemize}
    \item Distribution Mode Collapse: A collapse in the variety of samples generated for a single specific condition. The model fails to capture the probabilistic uncertainty of the future process.
    \item Mean Mode Collapse: A collapse in the distribution of the means of generated values across all different conditions. This effectively reduces the model to producing similar point estimates regardless of the historical input
\end{itemize}

When both occur simultaneously, 
the model yields a single point estimate for every input, 
completely losing its probabilistic forecasting capabilities, a scenario similar to outputs seen in the classical ForGAN experiments \ref{forgan_simple_sample_sigma}. 
Research indicates that mode collapse is often tied to the generator
converging toward sharp local minima. 
Once $G$ finds a sample that receives a high score from $D$,
it may continue to reproduce that sample. 
On the other hand,
if $D$ cannot identify this repetitive behavior (due to converging to its own local minima), 
the gradients for $G$ vanish, and the model becomes stuck.



The Fin-GAN paper \cite{forgan:fin_gan} identifies specific triggers and solutions for these performance issues caused by mode collapse in the ForGAN model. The key findings include:
\begin{itemize}
    \item Weight Initialisation: Using He initialisation led to mode collapse in 67\% of simulations for ForGAN; however, switching to Xavier initialisation significantly alleviated the problem.
    \item Loss Function Regularization: Including standard supervised terms like Mean Squared Error (MSE) or Sharpe Ratio (SR) with high coefficients can actually induce mode collapse by encouraging the model to only output values close to the target, so new terms are required to avoid this behaviour.
    \item The paper suggests that adding a Profit and Loss (PnL) term to the generator's loss function acts as an effective regularizer. This term shifts the generated distributions and prevents the model from collapsing into single-point estimates, ensuring a more robust probabilistic forecast.
\end{itemize}

Also in \cite{srforgan:srforgan}, the authors propose replacing the adversarial loss with a strictly proper scoring rule, specifically the Energy Score, which is designed to directly optimize the generator for accurate probabilistic forecasting.
This approach addresses the fundamental issues of mode collapse and training instability by providing a theoretically sound objective that encourages the generator to produce well-calibrated distributions rather than just realistic samples.\\

In particular, the paper proposes replacing the adversarial objective with a scoring rule whose theory is covered in \ref{sec:energy_score_theory}.
To sum up, it is a function $S(P^\phi, y)$ 
that takes a forecast distribution $P^\phi$ and a single observed realization $y$ and returns, in our framework, a penalty score where the lower the score, the better the forecast.\\
The key definition is strict propriety. A scoring rule $S$ is strictly proper relative to a class $\mathcal{P}$ if:
$$\mathbb{E}_{Y \sim P^*}[S(P^*, Y)] < \mathbb{E}_{Y \sim P^*}[S(P^\phi, Y)] \quad \forall P^\phi \neq P^*$$
Where the expected score is uniquely minimized when the forecast distribution equals the true data distribution $P^*$. No other distribution achieves a lower expected score. This is the property that makes a scoring rule the right objective for distribution learning because
it is theoretically guaranteed to converge to the right answer, provided that some conditions are met, not just some approximation like the 
GAN objective that minimizes a divergence between distributions which has a similar guarantee in principle, but the GAN only minimizes a variational lower bound of that divergence using a finite-capacity critic, which breaks the guarantee in practice.\\



%-----------------------------------------------------


The proposed model follows the scoring-rule-based conditional generative framework introduced by Pacchiardi et al. \cite{srforgan:srforgan}. Let $c\in\mathcal{C}$ denote the conditioning variable and let $y\in\mathcal{Y}$ denote the corresponding future realization. The objective is to learn a conditional probability distribution
\[
P^\phi(\cdot\mid c),
\]
parameterized by a neural network with weights $\phi$, such that it approximates the true conditional distribution
\[
P^*(\cdot\mid c).
\]

Unlike adversarial approaches, the model is trained by directly minimizing a strictly proper scoring rule. Given a scoring rule $S(P,y)$, the population objective is
\begin{equation}
\label{eq:population_objective}
\mathcal{L}(\phi)
=
\mathbb{E}_{(C,Y)\sim P^*}
\left[
S\!\left(P^\phi(\cdot\mid C),Y\right)
\right].
\end{equation}

In the conditional setting the strictly propriety of scoring rule is satisfied if:
\[
\mathbb{E}_{Y\sim P^*(\cdot\mid c)}
S(P^*(\cdot\mid c),Y)
<
\mathbb{E}_{Y\sim P^*(\cdot\mid c)}
S(P^\phi(\cdot\mid c),Y)
\]
for every conditional distribution
$
P^\phi(\cdot\mid c)\neq P^*(\cdot\mid c).
$
Consequently, the unique minimizer of (\ref{eq:population_objective}) satisfies
\[
P^\phi(\cdot\mid c)=P^*(\cdot\mid c)
\qquad
\text{for almost every }c,
\]
which provides the statistical consistency of the learning objective.\\


Before using any scoring rule as a training loss, \cite{srforgan:srforgan} and \cite{background:scoring_rules} require:

\begin{enumerate}
    \label{ES_conditions}
\item $\beta\in(0,2)$: the negative-power kernel $k(x,y)=-\|x-y\|^\beta$ is (conditionally) strictly positive-definite only in this range, which is what makes $S_E^{(\beta)}$ \ref{eq:energy_score} strictly proper rather than merely proper.
\item Finite $\beta$-th moments: $\mathbb{E}_{Y\sim P^\star}\|Y\|^\beta<\infty$ and $\mathbb{E}_{X\sim P^\phi}\|X\|^\beta<\infty$, without this the score, and hence the objective, is not even be finite.
\item A \emph{well-specified} unbiased estimator of the population scoring rule.
\item i.i.d.\ training pairs $(c_j,y_j)\sim\Pi\times P^\star(\cdot|c)$ 
\item A generator admitting the reparameterisation trick, $X=G_\phi(z,c)$, $z\sim\mathcal N(0,I)$, differentiable in $\phi$, so $\nabla_\phi S_E$ can be estimated unbiasedly (App.~C, \cite{srforgan:srforgan}).
\end{enumerate}
The discussion and proofs of these conditions are met in out setting are provided in the appendix \ref{ch:condition_proofs}.\\


The original formulation of Pacchiardi et al. \cite{srforgan:srforgan} considers a prequential forecasting problem in which the conditioning variables are overlapping sliding windows extracted from a single observed time series. 
In contrast, the present work constructs the training dataset by simulating $J$ independent trajectories of the underlying stochastic differential equation. 
Each trajectory produces one condition--target pair
\[
(c_j,y_j),
\qquad
j=1,\ldots,J,
\]
where
\[
c_j
=
(X_0^j,\ldots,X_N^j)
\]
is the observed trajectory up to the forecasting horizon and
\[
y_j
=
X_{N+\tau}^j
\]
is the future realization.



Although different trajectories may be generated using different model parameters (for example different values of the volatility parameter in the Black--Scholes model or different Heston parameters), this does not violate the theoretical assumptions underlying scoring-rule minimization. Indeed, if the parameters $\vartheta$ of the underlying data generation process are independently 
sampled from a probability distribution $\Pi(\vartheta)$ before simulating each trajectory, every observation is generated according to a hierarchical model. In fact, the first step of the generative process is to sample the parameters $\vartheta_j$ for each trajectory $j$ from $\Pi$, and then to sample the condition--target pair $(c_j,y_j)$ from the conditional law $P(C,Y \mid\vartheta)$ defined by those parameters. Formally:
\[
\vartheta_j\sim\Pi,
\qquad
(c_j,y_j)\sim P(C,Y \mid\vartheta),
\]
which induces the marginal joint distribution:
\[
(c_1,y_1),\ldots,(c_J,y_J)
\stackrel{\mathrm{i.i.d.}}{\sim}
P(C,Y),
\]

So, even though each trajectory may have been simulated with different underlying stochastic parameters. The variation in the parameters is absorbed into the marginal data distribution and does not affect the strict propriety of the scoring rule, which is defined conditionally on the observed conditioning variable.\\




The empirical objective optimized during training is therefore
\begin{equation}
\label{eq:empirical_sr}
\hat{\mathcal{L}}(\phi)
=
\frac{1}{J}
\sum_{j=1}^{J}
S_E\!\left(
P^\phi(\cdot\mid c_j),
y_j
\right),
\end{equation}
which is the standard empirical estimator of (\ref{eq:population_objective}) for independent condition--target observations.


%----------------------------------------------------



The paper \cite{srforgan:srforgan} uses the Energy Score as its primary scoring rule. 
For $\beta \in (0,2)$, with $X, X' \sim P^\phi$ independent draws:
\begin{equation}
    \label{eq:energy_score}
S_E^{(\beta)}(P^\phi, y) = 2 \cdot \mathbb{E}\|X - y\|_2^\beta - \mathbb{E}\|X - X'\|_2^\beta
\end{equation}

This Energy Score is obtained from the general Kernel Score:
\begin{equation*}
    \label{eq:kernel_score} 
    S_k(P^\phi, y) = \mathbb{E}[k(X,X')] - 2\mathbb{E}[k(X,y)]
\end{equation*}

By choosing the negative power kernel $k(x,y) = -\|x - y\|_2^\beta$. 
The Kernel Score is strictly proper whenever $k$ is a positive-definite kernel and $-\|x - y\|_2^\beta$ for $\beta \in (0,2)$ is indeed positive-definite.\\

The Energy Score can be decomposed in two terms:
$$S_E^{(\beta)}(P^\phi, y) = \underbrace{2\mathbb{E}\|X - y\|_2^\beta}_{\text{accuracy term}} - \underbrace{\mathbb{E}\|X - X'\|_2^\beta}_{\text{spread reward}}$$
The accuracy term pulls generated samples toward the truth. The spread reward is subtracted, meaning the optimizer is explicitly penalized for producing a degenerate distribution where $X \approx X'$. 
This is the mechanism that prevents mode collapse: a collapsed distribution has $\|X - X'\| = 0$, contributing nothing to the spread reward, making the total score worse than a properly dispersed distribution. The strictly proper guarantee ensures the only global minimum is at $P^\phi = P^*$.\\

The paper fixes $\beta = 1$ throughout. In 1D, $S_E^{(1)}$ recovers exactly the CRPS (Continuous Ranked Probability Score), which is the standard evaluation metric in probabilistic weather forecasting:
\begin{equation}
    \label{eq:crps}
    \begin{aligned}
        \mathrm{CRPS}(F, y) =& \int_{-\infty}^{\infty} \big(F(x) - \mathbf{1}\{x \ge y\}\big)^2 \, dx\\
        &\text{ or alternatively }\\
      \mathrm{CRPS}(F, y)  =& \mathbb{E}|X - y| - \frac{1}{2}\mathbb{E}|X - X'|
    \end{aligned}
\end{equation}


$S_E$ is defined as an expectation over $P^\phi$, but is not possible to differentiate through a sampling operation.
The reparameterization trick resolves this. Because the generator defines:
$$X = G_\phi(z, \text{c}), \quad z \sim \mathcal{N}(0, I)$$
the expectation formula can be rewritten as:
$$\mathbb{E}_{X \sim P^\phi}[\|X - y\|] = \mathbb{E}_{z \sim \mathcal{N}}[\|G_\phi(z, c) - y\|]$$
The randomness now lives entirely in $z$, which is fixed before the forward pass. 
The generator $G_\phi$ is deterministic and differentiable with respect to its weights. Automatic differentiation libraries can therefore compute $\nabla_\phi S_E$ exactly.
The unbiased Monte Carlo estimator of the full gradient, as discussed in Appendix C of the paper \cite{srforgan:srforgan} is:
\begin{equation}
    \label{eq:energy_score_estimator}
    \hat{S}_E^{(\beta)}(\{x_i\}_{i=1}^m, y) = \frac{2}{m}\sum_{i=1}^m \|x_i - y\|^\beta - \frac{1}{m(m-1)}\sum_{\substack{i,k=1 \\ k \neq i}}^m \|x_i - x_k\|^\beta
\end{equation}

where each $x_i = G_\phi(z_i, c)$ with $z_i$ drawn independently. The $\frac{1}{m(m-1)}$ denominator excludes diagonal terms and makes this an unbiased estimator. 
Moreover, the paper proves that exchanging gradient and expectation is valid almost surely under mild conditions on the neural network architecture. This means that minimizing the empirical energy score via SGD gives unbiased gradient estimates of the true objective, so convergence results from standard SGD theory apply.



\begin{figure}[H]
    \centering
    \includegraphics[width=1.0\textwidth]{SRForGAN/images/srforgan_train_schema.png}
    \caption{SRForGAN training schema. The generator takes a noise vector and a conditioning variable to produce a forecast distribution, which is then scored against the true observation using the Energy Score. The critic is not used in this architecture, as the adversarial loss is replaced by the scoring rule objective.}
    \label{img:srforgan_train_schema}
\end{figure}


%
%Layer 5 — The prequential formulation: adapting to time series
%This is the paper's main theoretical contribution. For a temporal sequence $(y_1, \dots, y_T)$, the standard i.i.d. scoring rule objective (Eq. 6 in the paper) is:
%$$\min_\phi \mathbb{E}_{\theta \sim \Pi} \mathbb{E}_{y \sim P^*(\cdot|\theta)} S(P^\phi(\cdot|\theta), y)$$
%This assumes independent context-observation pairs $(\theta_i, y_i)$, which is not true for a time series. The paper replaces this with the prequential (predictive-sequential) score. Given a recorded window of length $T$, forecast each step $l$ ahead conditioned on the previous $k$ observations:
%$$\hat{\phi}_T(y_{1:T}) := \arg\min_\phi \sum_{t=k}^{T-l} S(P^\phi_{t+l}(\cdot|y_{t-k+1:t}),\; y_{t+l})$$
%This is literally how forecasting systems are evaluated in practice: for each time window in your training data, generate a conditional forecast and score it against the realization. The prequential score is then the sum of these one-at-a-time forecasting evaluations across the entire training sequence.
%Why is this still a proper scoring rule despite temporal dependence? The key is Theorem 2 (in the paper): if $S$ is strictly proper, then the prequential score $S_T$ is strictly proper for distributions over $y_{k+l:T}|y_{1:k+l-1}$ that are $k$-Markovian with lag $l$. The proof proceeds by tower property of expectations — propriety at each step propagates to the joint proper scoring of the full sequence. Intuitively, because each term $S(P^\phi_{t+l}(\cdot|y_{t-k+1:t}), y_{t+l})$ depends only on the local window of length $k$, and the generator is the same for all $t$, the sequential evaluation correctly scores the generator's ability to forecast the conditional distribution of the next step.



%Layer 6 — The consistency theorem (Theorem 3)
%This is the rigorous guarantee that makes the method scientifically sound. The empirical minimizer:
%$$\hat{\phi}_T(y_{1:T}) = \arg\min_\phi \frac{1}{T-l-k+1}\sum_{t=k}^{T-l} S(P^\phi(\cdot|y_{t-k+1:t}), y_{t+l})$$
%converges to $\phi^*_T$, the minimizer of the true expected prequential score, as $T \to \infty$, under four assumptions:
%A1 — Compactness: $\Phi$ compact. Standard regularity condition.
%A2 — Uniqueness: The true minimizer is isolated. Lemma 12 shows this is satisfied when $S$ is strictly proper and the model is well-specified (i.e., there exists $\phi^*$ such that $P^{\phi^*} = P^*$).
%A3 — Asymptotic stationarity: The marginal distribution of $(y_{t-k+1:t+l})$ averaged over time converges weakly. This is satisfied for stationary processes like GBM, which has stationary increments.
%A4 — Mixing and moment boundedness: The process is either $\alpha$-mixing or $\phi$-mixing with appropriate decay rates. GBM satisfies $\phi$-mixing since increments are independent. The moment bound for the Energy Score reduces to requiring $\mathbb{E}\|y_{t+l}\|^\beta < \infty$, which holds for log-normal GBM increments.
%The proof combines two classical results: the Uniform Law of Large Numbers for dependent data (Pötscher and Prucha 1989) applied to the sum of scoring rule evaluations, and a generic consistency result (Theorem 5.1, Skouras 1998) which promotes uniform convergence of the objective to convergence of the minimizer. The key insight is that even though consecutive terms in the sum are correlated (since the windows overlap), the mixing condition ensures this dependence decays fast enough that the ULLN still holds. This is fundamentally different from GAN training, which has no analogous consistency result for dependent data.





% ---------------------------------------------------------------
% Algorithm 1 — Unbiased Energy Score Estimate
% ---------------------------------------------------------------
\begin{algorithm}[H]
\caption{Unbiased Energy Score Estimate $\hat{S}_E$}
\label{alg:energy_score}
\begin{algorithmic}[1]
\REQUIRE Generator $G_{\phi}$, condition $c \in \mathbb{R}^l$,
         observation $y \in \mathbb{R}^d$,
         exponent $\beta \in (0,2)$
\STATE Draw $m$ independent noise vectors $\{z_j\}_{j=1}^m \sim \mathcal{N}(0, I)$
\STATE Generate $m$ samples $\{x_j\}_{j=1}^m \subset \mathbb{R}^d$ where $x_j = G_{\phi}(z_j, c)$


\STATE \textbf{// Term 1: } $\frac{2}{m}\sum_{j=1}^{m}\|x_j - y\|^{\beta}$
\STATE $T_1 \leftarrow \dfrac{2}{m}\displaystyle\sum_{j=1}^{m}\|x_j - y\|^{\beta}$

\STATE \textbf{// Term 2: }
       $\frac{1}{m(m-1)}\sum_{j \neq k}\|x_j - x_k\|^{\beta}$
\STATE $T_2 \leftarrow \dfrac{1}{m(m-1)}\displaystyle\sum_{\substack{j,k=1 \\ j \neq k}}^{m}\|x_j - x_k\|^{\beta}$

\RETURN $\hat{S}_E \leftarrow T_1 - T_2$
\end{algorithmic}
\end{algorithm}

It is important to note that the available code implementation of the scoring rule includes also the 
kernel score, which is a generalization of the energy score. The kernel score is defined as:
\begin{equation}
    \label{eq:kernel_score_estimator}
    \hat{S}_k(\{x_i\}_{i=1}^m, y) = 
    \frac{1}{m(m-1)}\sum_{\substack{i,k=1 \\ k \neq i}}^m k(x_i, x_k) 
    - \frac{2}{m}\sum_{i=1}^m k(x_i, y), 
    \quad k(x,x') = \exp\!\left(-\frac{\|x - x'\|^2}{2\gamma^2}\right)
\end{equation}
where $\gamma$ is a bandwidth parameter. The kernel score is also strictly proper and can be used as an alternative to the energy score, but the paper focuses on the energy score for its simplicity and interpretability in terms of accuracy and spread.\\
One of the major proposition arising from \cite{srforgan:srforgan} is that the sum of the energy score and the kernel score is also a strictly proper scoring rule, which they call the energy-kernel score. This allows for a flexible combination of both scores in the generator's loss function, potentially leveraging the strengths of each to further improve forecasting performance.
The energy-kernel score is defined as:
\begin{equation}
    \label{eq:energy_kernel_score_estimator}
    \hat{S}_{E+k}^{(\beta)}(\{x_i\}_{i=1}^m, y) = 
    \hat{S}_E^{(\beta)}(\{x_i\}_{i=1}^m, y) 
    + \hat{S}_k(\{x_i\}_{i=1}^m, y)
\end{equation}
where $\hat{S}_E^{(\beta)}$ is the energy score estimator from Algorithm \ref{alg:energy_score} and $\hat{S}_k$ is the kernel score estimator from Equation \ref{eq:kernel_score_estimator}.


% ---------------------------------------------------------------
% Algorithm 2 — MySRForGAN Training
% ---------------------------------------------------------------
\begin{algorithm}[H]
\caption{\textsc{SRForGAN} Training via Energy Score Minimisation}
\label{alg:srforgan_train}
\begin{algorithmic}[1]
\REQUIRE
    Training set $\mathcal{D}_{\text{train}} = \{(y_{t+l},\, y_{t-k+1:t})\}_{t=k}^{T-l}$,
    validation set $\mathcal{D}_{\text{val}}$ (optional),
    generator $G_{\phi}$ with parameters $\phi$,
    scoring rule SR $\in \{\text{energy}, \text{kernel}, \text{energy-kernel}\}$,
    number of generator draws $m$,
    learning rate $\eta$,
    maximum epochs $E_{\max}$,
    early-stopping patience $P$,
    minimum improvement threshold $\delta_{\min}$
\ENSURE Trained generator parameters $\phi^*$

\STATE Initialise Adam optimiser with parameters $\phi$ and learning rate $\eta$
\STATE $\ell^*_{\text{val}} \leftarrow +\infty$;\quad
       $p \leftarrow 0$;\quad
       $\phi^* \leftarrow \phi$

\FOR{$e = 1, \ldots, E_{\max}$}

    \STATE \textbf{// --- Training pass ---}
    \STATE Shuffle $\mathcal{D}_{\text{train}}$ into mini-batches
           $\mathcal{B}_1, \ldots, \mathcal{B}_{N_B}$ of size $B$
    \STATE $\ell_{\text{train}} \leftarrow 0$

    \FOR{each mini-batch $(Y, C) \in \mathcal{B}_1, \ldots, \mathcal{B}_{N_B}$}
        \STATE $\hat{S}_E^B = 0$
        \FOR{$b = 1$ to $B$ samples in the mini-batch}
            \STATE \text{// Compute for sample $b$ of the batch  $\hat{S}_E^b $}
            \hfill\textit{// Algorithm~\ref{alg:energy_score}}
            \STATE $\hat{S}_E^B \leftarrow \hat{S}_E^B + \hat{S}_E^b$
        \ENDFOR
        \STATE $\hat{\mathcal{L}} \leftarrow
               \textsc{SRLoss}(C,\, Y,\, G_{\phi},) = \frac{1}{N_B} \hat{S}_E^B$
               \hfill\textit{// Estimator~\ref{eq:empirical_sr}}
        \STATE $\phi \leftarrow \phi - \eta\,\nabla_{\phi}\,\hat{\mathcal{L}}$
               \hfill\textit{// Adam update, no critic needed}
        \STATE $\ell_{\text{train}} \leftarrow \ell_{\text{train}} + \hat{\mathcal{L}}$
    \ENDFOR
    \STATE $\ell_{\text{train}} \leftarrow \ell_{\text{train}} / N_B$

    \STATE \textbf{// --- Validation pass (no gradient) ---}
    \IF{$\mathcal{D}_{\text{val}}$ is provided}
        \STATE $\ell_{\text{val}} \leftarrow \textsc{SRLoss}$ evaluated on
               $\mathcal{D}_{\text{val}}$ 
        \IF{$\ell_{\text{val}} < \ell^*_{\text{val}} - \delta_{\min}$}
            \STATE $\ell^*_{\text{val}} \leftarrow \ell_{\text{val}}$;\quad
                   $p \leftarrow 0$;\quad
                   $\phi^* \leftarrow \phi$
                   \hfill\textit{// save best checkpoint}
        \ELSE
            \STATE $p \leftarrow p + 1$
            \IF{$p \geq P$}
                \STATE \textbf{break}
                       \hfill\textit{// early stopping}
            \ENDIF
        \ENDIF
    \ENDIF

\ENDFOR

\RETURN $\phi^*$
\end{algorithmic}
\end{algorithm}





% ── local colour definitions ─────────────────────────────────────────────────
\definecolor{cganblue}{RGB}{30,100,180}
\definecolor{wganred}{RGB}{180,40,40}
\definecolor{commentgray}{RGB}{110,110,110}
 
% ── local math helpers ───────────────────────────────────────────────────────
\newcommand{\Loss}{\mathcal{L}}
\newcommand{\ES}{\hat{S}_{E}}
 
% ── line-comment for the algorithmic package ─────────────────────────────────
% algorithmic has no standalone comment line; emulate it with \STATE
\newcommand{\LINECOMMENT}[1]{%
  \STATE{\color{commentgray}\small$\triangleright$\;\textit{#1}}}
 
% override inline comment style
\renewcommand{\algorithmiccomment}[1]{%
  \hfill{\color{commentgray}\small$\triangleright$\;\textit{#1}}}
 
% ─────────────────────────────────────────────────────────────────────────────

Additionally, the paper states in the conclusion that the combination of adversarial and SR minimization would be beneficial and a promising path to explore. So an additional hybrid loss term is defined by adding the Energy Score to both the Wasserstein adversarial loss and classical Jensen-Shannon adversarial loss. 
In particular, the generator loss using the hybrid approach is defined as:
\begin{equation}
    \label{eq:hybrid_wasserstein_loss}
    \mathcal{L}_G = \underbrace{-\mathbb{E}[C(G(c,z))]}_{\text{Wasserstein signal}} + \lambda \cdot \underbrace{\hat{S}_E}_{\text{Energy Score}}
\end{equation}


 
The hybrid models, which use both the adversarial losses (Binary Cross-Entropy and Wasserstein distance) and the energy score, inherit the same energy score method from \ref{alg:energy_score}.\\
The hybrid Forgan uses as adversarial term the Binary Cross-Entropy loss and the
 full hybrid generator loss for one mini-batch $B$ is:
 
\begin{equation}
  \Loss_{G}^{\mathrm{CGAN}}
  =
  \underbrace{
    -\frac{1}{B} \sum_{b=1}^B \log D(G(z^{(b)}, c^{(b)}))
  }_{\text{BCE adversarial loss}}
  \;+\;
  \underbrace{
    \frac{\lambda}{B}\;
    \sum_{b=1}^B \ES\!\left(\{\hat{x}_m\}_{m=1}^{M},\,\mathbf{Y}\right)
  }_{\text{Energy Score (\ref{alg:energy_score})}}
  \label{eq:cgan_hybrid}
\end{equation}
 
\begin{algorithm}[H]
\caption{{\color{cganblue}\texttt{ForGAN}} --- hybrid generator update
  (called when $\mathrm{step} \bmod n_{\mathrm{critic}} = 0$)}
\begin{algorithmic}[1]
\REQUIRE Generator $G_{\phi}$, Discriminator $D_{\psi}$
\REQUIRE Condition batch $\mathbf{c}$, true targets $\mathbf{y}$,
         real labels $\mathbf{1}$, samples $M$,
         ES weight $\lambda$, noise dim $N_z$
\ENSURE  Updated generator parameters $\phi$


\LINECOMMENT{Draw M independent noise vectors; generate M samples}
\STATE $\mathrm{samples} \leftarrow [\;]$
\FOR{$m = 1$ \textbf{to} $M$}
  \STATE $\mathbf{z}_m \sim \mathcal{N}(\mathbf{0},\,\mathbf{I}_{N_z})$
  \COMMENT{fresh draw each iteration}
  \STATE $\hat{y}_m \leftarrow G_{\phi}(\mathbf{c},\,\mathbf{z}_m)$
  \COMMENT{shape $[B \times D]$}
  \STATE $\mathrm{samples}.\mathrm{append}(\hat{y}_m)$
\ENDFOR
 
\LINECOMMENT{Adversarial term: pass only the first sample through D}
\STATE $\mathrm{probs}
  \leftarrow D_{\psi}\!\left(\mathrm{samples}[1],\,\mathbf{c}\right)$
\STATE $\Loss_{\mathrm{adv}}
  \leftarrow \mathrm{BCE}(\mathrm{probs},\;\mathbf{1})$
  \COMMENT{generator wants $D$ to output 1 on fakes}
 
\LINECOMMENT{Energy Score term: all M samples contribute}
\STATE $\Loss_{\mathrm{ES}}
  \leftarrow
  \textsc{EnergyScoreLoss}(\mathrm{samples},\;\mathbf{y})$
 
\LINECOMMENT{Combine, backpropagate once, update}
\STATE $\Loss_{G}
  \leftarrow \Loss_{\mathrm{adv}} + \lambda \cdot \Loss_{\mathrm{ES}}$
\STATE $\mathrm{Backward}(\Loss_{G})$
  \COMMENT{gradients from both terms accumulate in one pass}
\STATE $\mathrm{Step}(\phi)$
\RETURN $\Loss_{G}$
\end{algorithmic}
\end{algorithm}


The adversarial term is the Wasserstein generator loss
$-\mathbb{E}[D(G(z, c))]$ so in this case no labels are needed since the critic
output is an unbounded real score.
 
\begin{equation}
  \Loss_{G}^{\mathrm{CWGAN}}
  =
  \underbrace{
    -\frac{1}{B} \sum_{i=1}^B D(G(z^{(b)}, c^{(b)}))}_{\text{Wasserstein generator loss}}
  \;+\;
  \underbrace{
    \frac{\lambda}{B}\;
    \sum_{i=1}^B \ES\!\left(\{\hat{x}_m\}_{m=1}^{M},\,\mathbf{Y}\right)
  }_{\text{Energy Score (\ref{alg:energy_score})}}
  \label{eq:cwgan_hybrid}
\end{equation}
 



\begin{algorithm}[H]
\caption{{\color{wganred}\texttt{ForWGAN}} --- hybrid generator update}
\begin{algorithmic}[1]
\REQUIRE Generator $G_{\phi}$, Critic $C_{\psi}$
\REQUIRE Condition batch $\mathbf{c}$, true targets $\mathbf{y}$,
         samples $M$, ES weight $\lambda$, noise dim $N_z$
\ENSURE  Updated generator parameters $\phi$
 
 
\LINECOMMENT{Draw M independent noise vectors; generate M samples}
\STATE $\mathrm{samples} \leftarrow [\;]$
\FOR{$m = 1$ \textbf{to} $M$}
  \STATE $\mathbf{z}_m \sim \mathcal{N}(\mathbf{0},\,\mathbf{I}_{N_z})$
  \STATE $\hat{y}_m \leftarrow G_{\phi}(\mathbf{c},\,\mathbf{z}_m)$
  \COMMENT{shape $[B \times D]$}
  \STATE $\mathrm{samples}.\mathrm{append}(\hat{y}_m)$
\ENDFOR
 
\LINECOMMENT{Wasserstein term: no labels; maximise critic score on fakes}
\STATE $\Loss_{\mathrm{W}}
  \leftarrow
  -C_{\psi}\!\left(\mathrm{samples}[1],\,\mathbf{c}\right).\mathrm{mean}()$
  \COMMENT{no target labels required}
 
\LINECOMMENT{Energy Score term: all M samples contribute}
\STATE $\Loss_{\mathrm{ES}}
  \leftarrow
  \textsc{EnergyScoreLoss}(\mathrm{samples},\;\mathbf{y})$
 
\LINECOMMENT{Combine, backpropagate once, update}
\STATE $\Loss_{G}
  \leftarrow \Loss_{\mathrm{W}} + \lambda \cdot \Loss_{\mathrm{ES}}$
\STATE $\mathrm{Backward}(\Loss_{G})$
\STATE $\mathrm{Step}(\phi)$
\RETURN $\Loss_{G}$
\end{algorithmic}
\end{algorithm}
 

The critic and the discriminator are trained on real and fake samples, respectively, and do not interact with the Energy Score in any way,
so they follow the standard adversarial training procedure already described in the previous chapter \ref{alg: cgan}, \ref{alg:cwgan_gp}.



\section{Performance metrics and statistical tests}
\label{sec:performance_metrics}
In the following section the performance metrics for probability distributions used to evaluate the proposed models will be described in details.
Moreover, the Kolmogorov-Smirnov test will be presented in order to evaluate the performance of the model in generating the future probability distribution of the asset price.\\
\subsection{Distance metrics}
When dealing with discretized probability distributions represented as vectors,
where $P = [p_1, p_2, \dots, p_k]$ and $Q = [q_1, q_2, \dots, q_k]$ are probability mass functions defined over the same binning, 
different metrics capture different geometric and information-theoretic properties.\\

The Total Variation (TV) distance is the $L_1$ distance between the vectors, defined as:
\begin{equation}
    \label{eq: TV_distance}
    \delta(P, Q) = \frac{1}{2} \sum |p_i - q_i|
\end{equation}
It represents the maximum possible difference between the probabilities that the two distributions assign to the same event,
but it is "bin-blind," meaning it doesn't care if the probability mass is moved to a neighboring bin or an outlier bin.\\
The Hellinger distance, defined as:
\begin{equation}
    H(P, Q) = \frac{1}{\sqrt{2}} \sqrt{\sum (\sqrt{p_i} - \sqrt{q_i})^2}
\end{equation}
It is a bounded metric ($[0, 1]$) which is more sensitive to differences in the "tails" of the distribution than TV-distance.\\
The Jensen-Shannon (JS) Divergence is a symmetric, smoothed version of the Kullback-Leibler divergence, calculated as:
\begin{equation}
    JS(P, Q) = \frac{1}{2}KL(P || M) + \frac{1}{2}KL(Q || M)
\end{equation}
where $M = \frac{1}{2}(P + Q)$. While JS is a popular metric for GANs, it suffers from the "disjoint support" problem where it reaches a constant maximum if the distributions do not overlap.\\
In contrast, the Wasserstein distance (specifically the 1D version for discretized vectors) can be computed as the $L_1$ distance between the Cumulative Distribution Functions (CDFs): 
\begin{equation}
    W_1(P, Q) = \sum |CDF_P(i) - CDF_Q(i)|
\end{equation}
Unlike the other three, the Wasserstein distance is "geometry-aware," meaning it incorporates the underlying distance between bins, making it far more robust for Monte Carlo simulations where we care about how far the generated values are shifted from the true values.

\subsection{The Kolmogorov-Smirnov test}
The Kolmogorov-Smirnov (KS) as defined in \cite{background:ks_test}, is a goodness of fit test that evaluates whether an empirical distribution is consistent with a fully specified continuous reference distribution. 
Let $F(x)$ denote the hypothesized cumulative distribution function (CDF), 
and let $F_{n}(x)$ be the empirical CDF constructed from a sample of size $n$. 
The KS statistic is defined as:
\begin{equation}
    D_{n}=\sup_{x}\lvert F_{n}(x)-F(x)\rvert
\end{equation}
representing the maximum vertical deviation between the empirical and theoretical CDFs.
Under the null hypothesis $H_{0}: F_{n}(x)=F(x)$,
the scaled distribution $\sqrt{n}D_{n}$ converges to the Kolmogorov distribution, enabling computation of p-values by using its CDF:
\begin{equation*}
    P(\sqrt{n}D_{n}\le t)=1-2\sum_{k=1}^{\infty}(-1)^{k-1}e^{-2k^{2}t^{2}}
\end{equation*}
Large values of $D_{n}$ lead to rejection of $H_{0}$, indicating
incompatibility between the empirical and theoretical distributions.\\

When applying the KS framework to generated outputs from the proposed GAN, the empirical object is a vector of bin probabilities
$(p_{1},\dots,p_{K})$.  These represent a discretized approximation of the probability distribution. To construct the empirical CDF, 
the cumulative probabilities are computed as:
\begin{equation*}
    F_{\text{GAN}}(x_{k})=\sum_{i=1}^{k}p_{i}
\end{equation*}

where $x_k$ indexes the bin boundaries. For comparison with a normal distribution, the theoretical CDF
$F_{\mathcal N}(x_{k})$ is computed at the same grid. The KS statistic becomes:
\begin{equation*}
    D=\max_{k}\lvert F_{\text{GAN}}(x_{k})-F_{\mathcal N}(x_{k})\rvert
\end{equation*}

This yields a valid discrepancy measure, although interpreting p-values requires specifying an effective sample size.
Nonetheless, $D$ remains a rigorous, distribution-free indicator of alignment between the GAN-generated distribution and the target normal CDF.
Actually, in order to compute the p\-value of the test, the survival function of the Kolmogorov distribution is used $P(\sqrt{n}D_n \geq \sqrt{n}\hat{D_n} | H_0)$ 
where $\hat{D_n}$ is the observed test statistic. 
The effective sample size $n$ has been set equal to the number of bins $B$ used to discretize the distribution.

